\documentclass[english,m,master,palatino,scrreprt,nogermanabstract]{WeSTthesis}
% Please read README.md for additional information on WeSTthesis parameters.

% Separate paragraphs by vertical space rather than by a first-line indent.
% The default indent-only style runs consecutive paragraphs together visually,
% which is hard to read in the long analytical passages of Chapters 6 and 7.
% parskip=half is KOMA's half-line setting, which separates paragraphs clearly
% without the loose look of a full blank line.
\KOMAoptions{parskip=half}

\usepackage[english,ngerman]{babel}      % English and new German spelling
\usepackage[utf8]{inputenc}              % correct input encoding
\usepackage[T1]{fontenc}                 % correct output encoding
\usepackage{graphicx}                    % enhanced support for graphics
\usepackage{tabularx}                    % more flexible tabular
\usepackage{booktabs}                    % professional-looking tables (\toprule etc.)
\usepackage{multirow}                    % multi-row table cells
\usepackage{amsfonts}                    % math fonts
\usepackage{amssymb}                     % math symbols
\usepackage{amsmath}                     % overall enhancements to math environment
\usepackage{caption}                     % caption formatting
\usepackage[numbers,sort&compress]{natbib} % \citep/\citet, numbered bracket style
\usepackage{microtype}                   % better justification, fewer overfull lines
\emergencystretch=3em                    % let long \texttt identifiers wrap rather than
                                          % overflow the margin (e.g. BradleyTerryRewardModel)
\usepackage[colorlinks=true, linkcolor=black, citecolor=black, urlcolor=black,
            pdfborder={0 0 0}]{hyperref} % clickable links, no colored boxes (load after natbib)
\usepackage{xurl}                        % allow \url to break at /, _, - (long code paths)
\urlstyle{tt}                            % keep \url in the same monospace style as \texttt
\usepackage{tikz}                        % pipeline diagram (Figure 5.1)
\usetikzlibrary{positioning,arrows.meta,shapes.geometric,calc,fit,backgrounds}
\usepackage{pgfplots}                    % charts
\pgfplotsset{compat=1.18}
\usepgfplotslibrary{groupplots}
% shared chart styling
\pgfplotsset{
  thesisplot/.style={
    width=0.86\textwidth, height=5.4cm,
    tick align=outside, tick pos=left,
    axis line style={gray!60}, grid=major,
    grid style={gray!20, line width=0.3pt},
    label style={font=\small}, tick label style={font=\footnotesize},
    legend style={font=\footnotesize, draw=gray!50, fill=white,
                  cells={anchor=west}},
  },
  holbar/.style={fill=blue!25, draw=blue!55!black, line width=0.4pt},
  gcabar/.style={fill=orange!45, draw=orange!65!black, line width=0.4pt},
}

% Colored callout boxes for highlighting key findings, definitions, and
% limitations, using the same restrained blue/orange/gray palette as the
% charts above rather than introducing new colors.
\usepackage[skins,breakable]{tcolorbox}
\newtcolorbox{keyfinding}[1][]{
  enhanced, colback=blue!6, colframe=blue!55!black,
  boxrule=0.5pt, left=2mm, right=2mm, top=1mm, bottom=1mm,
  fonttitle=\bfseries, coltitle=blue!30!black,
  attach boxed title to top left={yshift=-2mm, xshift=2mm},
  boxed title style={colback=white, boxrule=0pt, top=0.3mm, bottom=0.3mm},
  title={#1}, before skip=3mm, after skip=3mm,
}
\newtcolorbox{caveat}[1][]{
  enhanced, colback=orange!6, colframe=orange!65!black,
  boxrule=0.5pt, left=2mm, right=2mm, top=1mm, bottom=1mm,
  fonttitle=\bfseries, coltitle=orange!40!black,
  attach boxed title to top left={yshift=-2mm, xshift=2mm},
  boxed title style={colback=white, boxrule=0pt, top=0.3mm, bottom=0.3mm},
  title={#1}, before skip=3mm, after skip=3mm,
}
\newtcolorbox{definitionbox}[1][]{
  enhanced, colback=gray!6, colframe=gray!55!black,
  boxrule=0.5pt, left=2mm, right=2mm, top=1mm, bottom=1mm,
  fonttitle=\bfseries, coltitle=black,
  attach boxed title to top left={yshift=-2mm, xshift=2mm},
  boxed title style={colback=white, boxrule=0pt, top=0.3mm, bottom=0.3mm},
  title={#1}, before skip=3mm, after skip=3mm,
}

% The class typesets \theauthor in \Large on the title page; the second line drops
% back to \normalsize so the matriculation number sits under the name without
% competing with it.
\author{Muhammad Hasnat\\[.15cm]{\normalsize (Matriculation Number: 221100612)}}

\title{Benchmarking Holistic vs Sentence-Level Reward-Model Preferences for Factual Summarization via Granular Credit Assignment (GCA)}

\degreecourse{Web and Data Science}

\firstreviewer{Prof. Dr. Frank Hopfgartner}
\firstreviewerinfo{Institute for Web Science and Technologies}
\secondreviewer{Dr. Zeyd Boukhers}
\secondreviewerinfo{Fraunhofer Institute for Applied Information Technology FIT}
\mentor{Lingxiao Kong}
\mentorinfo{Fraunhofer Institute for Applied Information Technology FIT}

\englishabstract{%
\setlength{\parskip}{0em}%
Preference-based alignment of abstractive summarization models increasingly relies
on automatically generated preference labels rather than human judgments. It
remains unclear whether a single holistic judgment over a multi-sentence summary
uses the available supervision as effectively as a judgment assembled from
sentence-level evidence. This thesis studies that question as a controlled
comparison of two preference-construction procedures: a holistic condition, where
a complete candidate summary is scored against the source article, and a
Granular Credit Assignment (GCA) condition, where each summary is decomposed
into sentences, each is scored separately, and the sentence scores are
aggregated into a summary-level preference. The candidate summaries, the
factuality evaluator (AlignScore), the Bradley--Terry reward-model architecture,
and the training procedure are held fixed across four research questions, so
feedback granularity is isolated as the primary experimental variable.

\textbf{RQ1} asks whether GCA produces a more learnable preference signal, measured
as reward-model pairwise validation accuracy. Across twenty independently
seeded runs at $n=1{,}000$, GCA led in every
run, a mean advantage of 3.95 percentage points (bootstrap 95\% CI $[+3.19,
+4.70]$, run-level Wilcoxon $p<0.001$, Cohen's $d=2.22$). \textbf{RQ3} asks whether
this advantage is stable across dataset scale; it is not. Twenty-run campaigns at
$n=5{,}000$ and $n=10{,}000$, under the same corrected procedure, give mean
gaps of $-0.22$ and $-0.11$ percentage points, both statistically
indistinguishable from zero, with equivalence tests excluding any effect
larger than $\pm0.61$ and $\pm0.36$ points: the advantage is absent, not
merely undetected. Absolute accuracy improves for both conditions as data
grows even as their difference vanishes.

\textbf{RQ2} asks how evaluator configuration and aggregation strategy affect this
comparison. Both matter substantially: an early weakest-link penalty formula
hurt GCA relative to a simple mean, and a sweep across four AlignScore modes
finds GCA's advantage appears only under modes that do not themselves decompose
the input internally, suggesting external and evaluator-internal decomposition
act as partial substitutes. Each mode was tested once, so this is
associational, not confirmed.

\textbf{RQ4} asks whether the RQ1 advantage also holds against an independent,
held-out ground truth: 500 articles disjoint from every training set, scored by
the same judge but never used in training. It does, more sharply: at
$n=1{,}000$, GCA is significantly more precise than holistic on comparisons
with a genuine quality gap (up to $+9.25$ percentage points, $p<0.001$). It
does not merely fail to persist, as RQ1 does; it reverses to a significant
advantage for holistic at $n=10{,}000$ on three of four ground-truth
constructions ($p=0.010$--$0.036$). Two candidate explanations, near-zero-margin
training pairs and a length/sentence-count shortcut, were tested and both
rejected; the mechanism behind the reversal remains open.

These results characterize when granular feedback changes a preference
signal's structure, not that GCA is generally preferable: it is not, once
data are abundant or the comparison is ambiguous. Higher accuracy shows a
preference relation is more recoverable by this architecture, not that GCA
labels better match human judgments or that trained models are more
factually consistent, both future work.
}

\begin{document}
\pagenumbering{roman}
\maketitle

\selectlanguage{english}

% Front matter: acknowledgments, declaration on the use of AI assistance, and the
% list of abbreviations. Included from main.tex before the table of contents.

\chapter*{Acknowledgments}
\addcontentsline{toc}{chapter}{Acknowledgments}
\markboth{Acknowledgments}{Acknowledgments}

I would like to thank my first supervisor, Prof.\ Dr.\ Frank Hopfgartner of the
Institute for Web Science and Technologies, for his guidance throughout this
project and in particular for the feedback that led me to narrow the study to a
single, cleanly controlled comparison rather than a broader but less conclusive
pipeline. I am equally grateful to my second supervisor, Dr.\ Zeyd Boukhers of the
Fraunhofer Institute for Applied Information Technology FIT, for his technical
input on reward modeling and factuality evaluation. I owe particular thanks to
my mentor, Lingxiao Kong of the Fraunhofer Institute for Applied Information
Technology FIT, who guided me through this project from its early design
decisions to its final submission. His feedback shaped the direction of the
experiments, including the design of the independent ground-truth evaluation
in Chapter~\ref{ch:results}, and his close reading of multiple drafts of this
thesis directly produced the reorganization of the results around the research
questions rather than the chronology of the project, and much of the precision
in how the limitations are stated.

The experiments in this thesis were carried out on the MOGON NHR cluster at
Johannes Gutenberg University Mainz. I gratefully acknowledge the compute time made
available there; the repeated-run and multi-scale campaigns reported in
Chapter~\ref{ch:results} would not have been feasible without access to that
infrastructure.

Finally, I thank the University of Koblenz for providing the research environment
in which this work was completed, and my family and friends for their support over
the course of it.

\section*{Declaration on the Use of AI Assistance}

I acknowledge the use of large language models to support the writing of this
thesis, specifically for improving the clarity, structure, and grammatical
correctness of the text, and as a coding assistant during the implementation of the
experimental pipeline. The research design, the experiments, the analysis of the
results, and the conclusions drawn from them are my own. All generated text and
code were reviewed, verified, and revised by me, and I take full responsibility for
the content of this thesis.

\cleardoublepage

\tableofcontents
\cleardoublepage

\listoffigures
\cleardoublepage
\listoftables
\cleardoublepage

% List of abbreviations. Included from main.tex after the list of tables,
% so it sits with the other front-matter lists.

\chapter*{List of Abbreviations}
\addcontentsline{toc}{chapter}{List of Abbreviations}
\markboth{List of Abbreviations}{List of Abbreviations}

\begin{tabularx}{\textwidth}{@{}lX@{}}
\toprule
\textbf{Abbreviation} & \textbf{Full Term} \\
\midrule
AI               & Artificial Intelligence \\
CI               & Confidence Interval \\
CV               & Cross-Validation \\
DPO              & Direct Preference Optimization \\
GCA              & Granular Credit Assignment \\
GPU              & Graphics Processing Unit \\
LLM              & Large Language Model \\
MoDPO            & Multi-Objective Direct Preference Optimization \\
NHR              & Nationales Hochleistungsrechnen (National High-Performance Computing) \\
NLI              & Natural Language Inference \\
RLAIF            & Reinforcement Learning from AI Feedback \\
RLHF             & Reinforcement Learning from Human Feedback \\
RM               & Reward Model \\
ROUGE            & Recall-Oriented Understudy for Gisting Evaluation \\
\texttt{\_sp}    & Suffix marking the AlignScore evaluation modes that split the input
                   internally (\texttt{nli\_sp}, \texttt{bin\_sp}) \\
\bottomrule
\end{tabularx}

\cleardoublepage

\pagenumbering{arabic}

\chapter{Introduction}
\label{ch:introduction}

\section{Motivation}
\label{sec:intro-motivation}

Large language models (LLMs) are increasingly used for multi-sentence generation
tasks such as news summarization. In this setting, the perceived quality of an output is often
dominated not by its overall fluency but by a small number of high-impact factual
errors: a wrong entity, a flipped relation, or a fabricated claim can outweigh
otherwise well-written text. Throughout this thesis, \emph{reliability} refers
specifically to factual consistency with the source document, not to stylistic
polish or fluency.

Preference-based fine-tuning, most prominently reinforcement learning from human
feedback (RLHF), updates a model using comparative judgments between candidate
outputs rather than token-level supervision \citep{christiano2017deep,
stiennon2020learning, ouyang2022training}. Collecting high-quality human preferences
at scale is expensive and slow, so reinforcement learning from AI feedback (RLAIF)
replaces human annotators with an automatic judge \citep{bai2022constitutional,
lee2023rlaif}. This trades annotation cost for a new source of label noise and bias.

A further question arises for long, multi-sentence outputs such as news summaries,
one that is largely independent of whether the feedback comes from a human or an AI
judge: how granular should the preference signal be? A single holistic judgment
(``summary A is better than summary B'') collapses all of the evidence in a
multi-sentence output into one binary decision. If a summary is mostly accurate but
contains one fabricated sentence, a holistic judge may still prefer it over an
alternative that is uniformly mediocre but contains no outright fabrication. The
holistic label cannot express that distinction. Sentence-level judging can in
principle localize exactly which claims are, and are not, supported by the source,
and that localized evidence can then be aggregated into a summary-level decision.
This thesis asks whether that extra granularity actually produces a better training
signal, and under what conditions.

\section{Problem Statement}
\label{sec:intro-problem}

When constructing AI-generated preference data for training a factuality-oriented
reward model, it is unclear whether coarse, holistic preferences are sufficient, or
whether they systematically underuse the supervision signal available in long,
multi-sentence outputs. The problem addressed here, under controlled conditions, is
whether sentence-level factuality judgments, aggregated into a summary-level
preference through a mechanism referred to as \emph{Granular Credit Assignment}
(GCA), yield a more learnable preference signal for reward-model training than
purely holistic judgments of the same candidate summaries.

\begin{definitionbox}[Learnability, as used throughout this thesis]
The pairwise validation accuracy of a Bradley--Terry reward model trained on
the resulting preference pairs (Chapter~\ref{ch:background},
Section~\ref{sec:bg-bradley-terry}): given a held-out preference pair, how
often does the trained reward model rank the summary the judge preferred
above the alternative?
\end{definitionbox}

\section{Thesis Objective}
\label{sec:intro-objective}

The objective of this thesis is to determine, through a controlled empirical
comparison, whether Granular Credit Assignment produces a more learnable
factuality-preference signal than holistic scoring. The two conditions differ only
in how factuality scores are obtained and converted into a pairwise label:

\begin{quote}
\textbf{Holistic:} article $+$ complete candidate summary $\rightarrow$ factuality
score $\rightarrow$ pairwise preference.

\textbf{GCA:} article $+$ individual sentences $\rightarrow$ sentence factuality
scores $\rightarrow$ aggregation $\rightarrow$ pairwise preference.
\end{quote}

Both preference sets are then used to train the same Bradley--Terry reward-model
architecture under the same procedure. Holding the source articles, the candidate
summaries and their generation settings, the evaluator family, the reward-model
architecture, and the training procedure fixed allows feedback granularity to be
studied as the primary experimental variable, rather than as one difference among
several.

The thesis does not propose a new reward-model architecture, a new factuality
metric, or a new training objective. Its subject is the construction of the
preference data itself.

\section{Research Questions}
\label{sec:intro-rqs}

Four research questions follow from this objective. They are stated here in
summary form; Chapter~\ref{ch:methodology},
Section~\ref{sec:methodology-rqs} gives them together with the working hypotheses
and the scope conditions under which each can be answered, and
Chapter~\ref{ch:conclusion}, Section~\ref{sec:conclusion-rqs} answers them
against the evidence collected.

\begin{itemize}
  \item \textbf{RQ1.} Does sentence-level factuality feedback aggregated via GCA
    produce a more learnable pairwise preference signal than holistic
    summary-level feedback, as measured by reward-model pairwise validation
    accuracy?

  \item \textbf{RQ2.} How do evaluator configuration and sentence-score
    aggregation strategy affect the relative learnability of holistic and
    GCA-derived preferences?

  \item \textbf{RQ3.} How stable are the observed differences across repeated
    reward-model training runs and dataset sizes?

  \item \textbf{RQ4.} Does the learnability advantage measured in RQ1 also hold
    when the trained reward models are evaluated against an independent,
    held-out ground truth, rather than against their own training-style
    labels, and does that hold as the training-set size grows?
\end{itemize}

RQ1 is the confirmatory question and is addressed by a repeated-run comparison.
RQ2 is addressed by an exploratory sweep, and therefore supports statements about
association between configuration and outcome rather than about causal mechanism.
RQ3 motivates both the repeated runs at a single dataset size and the
multi-scale check at $n=5{,}000$ and $n=10{,}000$. RQ4 was added after that
multi-scale check, to test the learnability finding against data the trained
models never saw and never influenced their selection
(Chapter~\ref{ch:methodology}, Section~\ref{sec:methodology-ground-truth}).

\section{Contributions}
\label{sec:intro-contributions}

This work makes the following contributions:

\begin{enumerate}
  \item A controlled framework for comparing holistic and sentence-level,
    GCA-based construction of factuality preferences for abstractive
    summarization, in which the candidate pool, evaluator, reward-model
    architecture, and training procedure are held fixed
    (Chapter~\ref{ch:methodology}).

  \item An empirical study of how evaluator mode and aggregation strategy
    interact with feedback granularity. Neither the aggregation formula nor the
    evaluator configuration is neutral: a penalty-weighted aggregation performs
    worse than a simple mean, and in an exploratory sweep the evaluator mode is
    associated with whether an advantage for GCA is observed at all
    (Chapter~\ref{ch:results}, Sections~\ref{sec:results-formula-opt}
    and~\ref{sec:results-mode-sweep}).

  \item A repeated-run, multi-scale evaluation establishing a statistically
    significant GCA advantage at $n=1{,}000$
    (20 independent runs, mean $+3.95$ percentage points, run-level
    Wilcoxon $p<0.001$; found during final review to have been trained
    under AlignScore's \texttt{nli\_sp} mode rather than the \texttt{nli}
    mode intended, Chapter~\ref{ch:implementation},
    Section~\ref{sec:impl-preferences}), and confirming, through twenty-run campaigns at
    $n=5{,}000$ and $n=10{,}000$ matching that same resolution, that this
    advantage is genuinely absent
    rather than merely unreproduced at larger scale: both confidence intervals
    are centred on zero and neither overlaps the $n=1{,}000$ result
    (Sections~\ref{sec:results-final-campaign} and~\ref{sec:results-scale}).

  \item An exploratory observation that the benefit of external sentence-level
    decomposition may be reduced when the underlying evaluator already performs
    comparable internal decomposition, offered as a hypothesis motivated by the
    mode comparison rather than as an established mechanism
    (Chapter~\ref{ch:discussion}, Section~\ref{sec:disc-interpretation}).

  \item An independent, held-out ground-truth evaluation (RQ4) showing that the
    $n=1{,}000$ advantage is not an artefact of validating against the trained
    models' own training-style labels: it holds, more sharply, against 500
    articles disjoint from every training set used in this thesis, and it
    reverses rather than merely vanishing as the training set grows, becoming a
    significant advantage for holistic scoring at $n=10{,}000$
    (Chapter~\ref{ch:results}, Section~\ref{sec:results-ground-truth}). Two
    candidate explanations for the reversal were tested directly and both
    rejected (Sections~\ref{sec:results-gt-margin}
    and~\ref{sec:results-gt-shortcut}).

  \item A transparent account of instability, stochastic variation across runs,
    methodological limitations, and the reproducibility of the pipeline
    (Chapters~\ref{ch:implementation}, \ref{ch:discussion},
    and~\ref{ch:conclusion}).
\end{enumerate}

These contributions concern the construction and learnability of preference data.
The thesis does not claim that GCA yields objectively more accurate factuality
labels, nor that it improves the factual consistency of generated summaries; the
scope of what the experiments support is stated precisely in
Chapter~\ref{ch:methodology}, Section~\ref{sec:methodology-scope}.

\section{Scope}
\label{sec:intro-scope}

\begin{itemize}
  \item \textbf{Task:} single-document abstractive news summarization, using the
    CNN/DailyMail dataset \citep{hermann2015teaching, nallapati2016abstractive}.
  \item \textbf{Candidate generation:} a single instruction-tuned base model,
    Mistral-7B-Instruct-v0.3, sampled at two temperatures ($T=0.7$ and $T=1.0$) to
    produce an A/B candidate pair per source article.
  \item \textbf{Judge:} a single fixed, deterministic factuality judge,
    \texttt{yzha/AlignScore} \citep{zha2023alignscore}, intended throughout
    in its \texttt{nli} evaluation mode, though the confirmatory campaigns
    were, due to an implementation oversight found during final review,
    actually run under its \texttt{nli\_sp} mode
    (Chapter~\ref{ch:implementation}, Section~\ref{sec:impl-preferences}).
    No generative LLM-as-judge and no external API dependency is
    used in the final pipeline.
  \item \textbf{Optimization target:} the learnability of the resulting preference
    data for Bradley--Terry reward-model training, measured by pairwise validation
    accuracy under 5-fold cross-validation. A downstream fine-tuned generation
    policy is not evaluated in the final comparison; an early Direct Preference
    Optimization (DPO) prototype was explored but dropped from the final scope
    (Chapter~\ref{ch:methodology}).
\end{itemize}

\section{Thesis Structure}
\label{sec:intro-structure}

Chapter~\ref{ch:related-work} situates this work within the literature on
preference-based alignment, LLM-as-a-judge reliability, fine-grained feedback for
long-form generation, and automatic factual-consistency evaluation.
Chapter~\ref{ch:background} introduces the technical mechanisms this thesis builds
on: the Bradley--Terry preference model, Direct Preference Optimization, the
AlignScore factuality metric, and the statistical tools
used to validate the results. Chapter~\ref{ch:methodology} presents the experimental
design, including the GCA aggregation formula itself, which is this thesis's own
contribution rather than prior work, and an explicit account of how the study that
was actually carried out diverges from the original proposal, and why. Chapter~\ref{ch:implementation}
describes the pipeline, codebase, and compute infrastructure. Chapter~\ref{ch:results}
reports the full experimental campaign, from the initial reward-model baseline
through the formula and judge-mode optimization to the final multi-seed and
multi-scale validation, and then to the independent ground-truth precision
evaluation that answers RQ4. Chapter~\ref{ch:discussion} interprets these
results, including why the ground-truth result reverses rather than merely
fading where the learnability result only fades, and
Chapter~\ref{ch:conclusion} answers the research questions, states the
limitations of this work, and outlines directions for future work.

\chapter{Related Work}
\label{ch:related-work}

This chapter reviews four bodies of work that this thesis draws on: preference-based
alignment and how preference data becomes a training signal
(Section~\ref{sec:rw-rlhf-rlaif}); the reliability of AI judges when they replace
human annotators (Section~\ref{sec:rw-judge-reliability}); the granularity of
feedback in long-form generation, which is the variable this thesis manipulates
(Section~\ref{sec:rw-finegrained}); and automatic factual-consistency evaluation,
which supplies the judge (Section~\ref{sec:rw-factual-eval}).
Section~\ref{sec:rw-synthesis} draws these together into the specific gap this
thesis addresses, and Section~\ref{sec:rw-positioning} states what separates the
present work from the studies closest to it.

\section{Preference-Based Alignment: RLHF, RLAIF, and Offline Optimization}
\label{sec:rw-rlhf-rlaif}

Preference-based alignment collects pairwise comparisons between candidate model
outputs and converts those comparisons into a training signal. The approach
originates in preference learning for reinforcement learning, where human
annotators compare short trajectory segments and a reward model is fitted to
reproduce their comparisons, allowing an agent to be trained on objectives that
are easy to judge but hard to specify numerically \citep{christiano2017deep}. The
choice to elicit comparisons rather than absolute scores is deliberate: relative
judgments between two concrete outputs are more consistent between annotators than
absolute ratings on an abstract scale, and the Bradley--Terry model
(Chapter~\ref{ch:background}, Section~\ref{sec:bg-bradley-terry}) provides a
principled way to recover a scalar reward from them.

\citet{stiennon2020learning} adapted this recipe to summarization, and their
pipeline is the direct ancestor of the one used in this thesis. They collect human
comparisons between candidate summaries, train a reward model on those comparisons,
and then fine-tune a policy against the learned reward with reinforcement learning.
Two of their findings matter here. First, the reward model's own held-out agreement
with human preferences is reported as a diagnostic in its own right, which
establishes the precedent for treating reward-model accuracy as a meaningful
measurement rather than merely an intermediate artifact, as this thesis does.
Second, they observe that optimizing too hard against a learned reward degrades
true quality, an early instance of the reward-overoptimization problem that
motivates careful attention to the quality of the preference signal itself. RLHF
was subsequently scaled to general instruction-following \citep{ouyang2022training},
establishing the reward-model-plus-policy-optimization pattern as the default
alignment recipe.

Because human comparison data is expensive, reinforcement learning from AI feedback
(RLAIF) substitutes a model for the human annotator. Constitutional AI generates
critiques and revisions from a written set of principles and derives preference
labels from them, removing human annotation from the harmlessness training loop
almost entirely \citep{bai2022constitutional}. \citet{lee2023rlaif} compare RLAIF
and RLHF directly on summarization and dialogue tasks and report broadly
comparable downstream performance, which is the empirical result that makes
AI-generated preference data a credible substitute rather than a compromise. They
also emphasize that the outcome depends substantially on how the judge is
prompted and calibrated, a dependency that recurs throughout
Section~\ref{sec:rw-judge-reliability} and, in a different form, in the
evaluator-configuration effects observed here
(Chapter~\ref{ch:results}, Section~\ref{sec:results-mode-sweep}).

On the optimization side, RLHF pipelines have typically used on-policy algorithms
such as Proximal Policy Optimization \citep{schulman2017ppo}, which require
sampling from the policy during training and maintaining a separate reward model.
Direct Preference Optimization (DPO) removes both requirements by reparameterizing
the RLHF objective so that the optimal policy can be recovered in closed form
directly from preference pairs \citep{rafailov2023dpo}. DPO was the objective
originally proposed for this thesis, chosen because it isolates the effect of the
preference data from the dynamics of reinforcement learning. As
Chapter~\ref{ch:methodology} describes, DPO was implemented and produced early
results before being dropped from the final comparison in favor of evaluating
Bradley--Terry reward models directly. That change narrows the comparison further
than DPO alone would: it removes both the optimization dynamics that motivated
adopting DPO and the decoding variance introduced by generating and scoring a
fine-tuned policy's outputs, leaving a more targeted question about whether the
preference data itself is more learnable.

\section{Reliability Risks in AI Supervision and LLM-as-a-Judge}
\label{sec:rw-judge-reliability}

Substituting a language model for a human annotator introduces a distinct class of
error. Benchmark studies of LLM-as-a-judge document sensitivity to prompt wording,
position bias favoring whichever candidate appears first, verbosity bias favoring
longer answers regardless of quality, inconsistent scoring across repeated calls,
and self-enhancement effects in which a judge prefers outputs from models similar
to itself \citep{zheng2023judging}. \citet{liu2023geval} document related
calibration problems in LLM-based evaluation of generated text, and a recent
survey synthesizes the accumulated evidence on these failure modes and the
mitigations proposed for them, including rubric enforcement, pairwise rather than
absolute evaluation, and position swapping \citep{gu2024survey}.

These issues are more consequential in RLAIF than in evaluation. An unreliable
evaluation metric produces a noisy estimate of a model's quality; an unreliable
judge used to construct training data produces systematically mislabeled
supervision, and the model is then optimized to fit that mislabeling.
\citet{casper2023open} catalog failure modes of this kind across the RLHF
pipeline, including reward misspecification and the divergence between training
loss and intended behavior, and argue that problems in the preference data
propagate into the trained policy in ways that are difficult to detect after the
fact.

One response to this literature is to treat the judge as a fallible instrument and
wrap it in a reliability protocol: confidence gating, order randomization,
evidence-grounded rationale requirements, and prompt-robustness checks. The design
adopted here removes the exposure at its source instead, by using a fixed,
deterministic factuality model rather than a generative judge
(Section~\ref{sec:rw-factual-eval}). A deterministic model cannot exhibit prompt
sensitivity, position bias, or run-to-run inconsistency, because it is never
prompted and returns no natural-language output. This eliminates the category of
risk the protocol above is designed to manage, at two costs: it does not make the
evaluator correct, and it removes the rationale text that a human audit would
examine. Chapter~\ref{ch:methodology}, Section~\ref{sec:methodology-scope} states
what this implies for the claims the thesis can support.

\section{Feedback Granularity and Credit Assignment in Long-Form Generation}
\label{sec:rw-finegrained}

The variable this thesis manipulates is how localized the supervision signal is.
A single preference over a multi-sentence output provides one bit of information
about a text that may contain a dozen independently checkable claims, and it does
not indicate which claim caused the preference. This is a credit assignment
problem, and it is present even in the original preference-learning formulation:
\citet{christiano2017deep} have annotators compare short trajectory segments
rather than whole episodes precisely to make the resulting supervision denser and
easier to attribute.

The closest prior work to this thesis is \citet{wu2023finegrained}, who study
fine-grained human feedback for long-form question answering. Their approach
differs from holistic RLHF along two dimensions simultaneously. First, feedback is
collected at the density of individual sentences or sub-sentence spans rather than
whole responses, so an annotator marks the specific span that is irrelevant,
untruthful, or incomplete. Second, they train several distinct reward models, one
per error category, and combine their outputs into a dense reward signal during
reinforcement learning, so the policy receives a reward at each segment rather than
a single scalar at the end of the sequence. They report that this improves both the
resulting policy and the reward models themselves relative to holistic feedback,
and that the multiple reward models allow the relative weight of different
objectives to be adjusted after training.

Three differences separate their setting from this thesis. Their fine-grained
labels come from trained human annotators, whereas this thesis asks whether the
same localization strategy is still useful when the localized judgments come from
an automatic judge, which is cheaper by orders of magnitude but noisier per label.
Their fine-grained signal is consumed as a dense per-segment reward during
reinforcement learning, whereas GCA aggregates local scores back into a single
summary-level preference pair so that the resulting dataset remains compatible
with standard pairwise preference optimization, requiring no change to the training
objective. And they vary feedback density and objective decomposition together,
whereas this thesis holds the objective fixed at factual consistency alone and
varies only granularity, which isolates the granularity question at the cost of
addressing a narrower one.

Three more recent lines of work place sentence- or step-level signals at
other points in the alignment pipeline, and each should be distinguished
carefully from the present study, since none leaves the question examined
here open in general.

\citet{gu2025maskdpo} introduce Mask-DPO, which incorporates sentence-level
factuality signals directly into the DPO objective. Rather than treating a whole
response as preferred or dispreferred, Mask-DPO uses sentence-level factuality
information to learn only from the factually correct sentences within a preferred
response, and to avoid penalizing factually correct content that appears inside a
dispreferred one. The motivation is that response-level factuality supervision is
noisy, because preferred and dispreferred responses each mix accurate and
inaccurate sentences. The intervention point is the policy-optimization objective:
the preference pair remains a response-level object, and the granular signal is applied as a mask during training. The present thesis intervenes earlier and
leaves the training objective untouched. It asks how automatic factuality
judgments at different granularities should be converted into the pairwise label
in the first place, and measures the learnability of the resulting labels under a
fixed Bradley--Terry architecture.

\citet{qiu2025sentencelevel} modify the reward model itself, proposing a
sentence-level reward model that assigns scores at sentence boundaries and
aggregates them into a response-level score, while still learning from
response-level human preferences. Their aim is to address reward sparsity by
producing an intermediate granularity between response-level and token-level
reward modeling. The distinction from this thesis is architectural: they change
what the reward model computes, and the preference labels they learn from are
conventional response-level human judgments. This thesis does not modify the
reward-model architecture at all. The Bradley--Terry model, its backbone, and its
training procedure are held fixed across both conditions precisely so that
differences in outcome can be attributed to the preference labels rather than to
the model consuming them.

Closer still to GCA's decompose-then-aggregate structure is process reward
modeling for multi-step reasoning \citep{lightman2023verify}. Rather than
training a single model to judge a complete solution as correct or
incorrect, \citet{lightman2023verify} have human annotators label the
correctness of each individual reasoning step in a chain-of-thought
solution, train a reward model to reproduce those per-step judgments, and
show that this process-level supervision outperforms a response-level
outcome reward for selecting correct solutions among many sampled
candidates. The structural resemblance to GCA is real: both decompose a
long output into smaller units, score each unit, and use the per-unit
scores rather than a single holistic judgment. Two differences remain.
Process reward models keep the per-step scores as the object consumed
downstream, typically for reranking or search over sampled solutions,
whereas GCA aggregates its per-sentence scores back into one summary-level
score so that the output is a conventional pairwise preference, compatible
with unmodified Bradley--Terry training. And the per-step labels in
\citet{lightman2023verify} come from trained human annotators judging
mathematical correctness, a domain with a comparatively well-defined notion
of a correct step, whereas GCA's per-sentence labels come from an automatic
factuality judge applied to open-ended news summaries, where what counts as
a faithful sentence is judged by a model rather than verified against a
solution. The result that step-level supervision outperforms outcome-level
supervision in their setting is consistent with the motivation for GCA, but
it is evidence from a different domain, a different label source, and a
different point of consumption in the pipeline, not a direct precedent for
the automatic, aggregated case this thesis studies.

Reward models are also increasingly evaluated in their own right rather than only
as components of a larger pipeline. RewardBench assembles prompt-chosen-rejected
trios across several domains and evaluates reward models by how often they rank
the preferred response above the dispreferred one
\citep{lambert2024rewardbench}. That pairwise-ranking accuracy is the same
quantity used as the dependent variable in this thesis
(Chapter~\ref{ch:methodology}, Section~\ref{sec:methodology-scope}), which places
the measurement on established ground; the difference is that RewardBench varies
the reward model against fixed data, whereas this thesis varies the data
construction against a fixed reward model.

\section{Factual Consistency in Summarization and Evaluation}
\label{sec:rw-factual-eval}

News summarization remains a standard benchmark domain, with CNN/DailyMail widely
used for both training and evaluation \citep{hermann2015teaching,
nallapati2016abstractive}. Abstractive systems have long produced fluent output
that nonetheless contradicts or embellishes the source \citep{see2017get}, and this
persistent gap between fluency and faithfulness is what makes factual consistency
a worthwhile optimization target rather than a solved preprocessing concern.

Automatic factuality evaluation has developed along two main lines. Entailment-based
methods treat the source as a premise and the summary as a hypothesis: FactCC
trains a classifier on synthetically corrupted summaries to detect factual
conflicts \citep{kryscinski2020evaluating}. Question-answering-based methods
instead generate questions from the summary and check whether the source yields
the same answers, an approach developed in FEQA \citep{durmus2020feqa} and refined
in QAFactEval, which systematically compares components of the QA-based pipeline
\citep{fabbri2022qafacteval}.

\paragraph{Decomposition and aggregation in SummaC.} The work most directly
relevant to GCA's design is SummaC \citep{laban2022summac}, which observes that
natural language inference models trained on sentence-level data transfer poorly
to document-level consistency checking, partly because the inputs are far longer
than anything seen in training. Their solution is to decompose the problem: they
compute an entailment matrix between each source sentence and each summary
sentence, then aggregate that matrix into a single score. Critically for this
thesis, they show that \emph{how} the matrix is aggregated changes performance
substantially. Their zero-shot variant takes, for each summary sentence, the
maximum entailment score across all source sentences, and then averages these
maxima across summary sentences. Their learned variant instead trains a
convolutional layer over the distribution of entailment scores per summary
sentence, retaining information that the maximum discards, and performs better.

That aggregation strategy is a design decision with measurable consequences is
therefore established in the evaluation literature before this thesis takes it up.
The contribution here is to ask the same question one step upstream, at the point
where scores are converted into training labels rather than into an evaluation
metric, and Chapter~\ref{ch:results}, Section~\ref{sec:results-formula-opt} finds
the same lesson holds: the choice between a mean and a penalty-weighted
aggregation changes the outcome more than the decision to decompose at all.

\paragraph{AlignScore.} The judge used throughout this thesis, AlignScore, was not
part of the literature surveyed in the original proposal. \citet{zha2023alignscore}
recast a broad collection of natural language processing tasks, including natural
language inference, question answering, paraphrase detection, and fact
verification, into a single ``information alignment'' format, and train one model
on the union. The resulting function takes a context and a claim and returns a
scalar in $[0,1]$ estimating whether the context supports the claim. Because it is
trained across many task formats rather than on entailment data alone, it
generalizes across factuality benchmarks better than the single-task methods above.

AlignScore also inherits the decomposition question from SummaC, and resolves it
internally: in its split evaluation modes, the context is divided into chunks and
the claim into sentences, and the final score aggregates the resulting matrix by
taking a maximum over chunks and then a mean over claim sentences, which is
structurally the same rule as SummaC's zero-shot variant. This internal behavior
turns out to matter a great deal for the present work.
Chapter~\ref{ch:results}, Section~\ref{sec:results-mode-sweep} finds that whether
GCA outperforms holistic scoring depends on which AlignScore mode is used, and
Chapter~\ref{ch:discussion}, Section~\ref{sec:disc-interpretation} argues that this
is because the split modes already perform, inside the judge, much of the
decomposition that GCA performs outside it.

\paragraph{Fine-grained evaluation.} More recent work evaluates summaries at
sub-summary granularity rather than assigning a single score. FineSurE assesses
faithfulness at the sentence level and content coverage against a list of extracted
key facts, producing a diagnostic profile rather than one number
\citep{song2024finesure}. Like SummaC, and unlike this thesis, it is positioned as
an evaluation framework: it tells a practitioner where a summary went wrong, but
does not convert that localized judgment into supervision for training.

FActScore decomposes a generation into individual atomic facts and checks
each against a knowledge source, then reports the fraction supported as a
single precision score \citep{min2023factscore}. Its decompose-score-aggregate
structure is the closest evaluation-side analogue to GCA's own aggregation
step (Chapter~\ref{ch:methodology}, Section~\ref{sec:bg-gca}): both split a
generation into small checkable units, score each independently, and combine
the results into one number. The differences are the same ones that
separate this thesis from SummaC and FineSurE. FActScore decomposes into
atomic facts extracted from the generation itself and verifies them against
an external knowledge source, whereas GCA decomposes into sentences and
scores each against the source article with a single fixed judge; and
FActScore's aggregate score is reported as a diagnostic metric, not
converted into a pairwise preference for training a reward model. The
aggregation choice it makes, a simple fraction of atomic facts supported,
is nonetheless a further data point for the same lesson this thesis draws
in Section~\ref{sec:results-formula-opt}: how per-unit scores are combined
is a design decision with measurable consequences, examined here independently
by two unrelated methods on two different granularities of decomposition.

\section{Synthesis: The Gap This Thesis Addresses}
\label{sec:rw-synthesis}

Table~\ref{tab:rw-gap} summarizes the closest lines of work along the dimension
that distinguishes the present study: the point in the pipeline at which
sentence-level information is introduced, and what each approach holds fixed while
doing so.

\begin{table}[htbp]
\centering
\caption[Where sentence-level information enters the pipeline in related work]{Where sentence-level information enters the pipeline in related work.}
\label{tab:rw-gap}
\begin{tabularx}{\textwidth}{lXX}
\toprule
Work & Sentence-level signal enters at & Held fixed \\
\midrule
\citet{stiennon2020learning} & not used & -- \\
\citet{lee2023rlaif}         & not used & -- \\
\citet{wu2023finegrained}    & reward signal during RL (human labels) &
  response-level objective \\
\citet{gu2025maskdpo}        & the DPO objective, as a sentence mask &
  preference pair construction \\
\citet{qiu2025sentencelevel} & the reward-model architecture &
  preference labels (human, response-level) \\
\citet{lightman2023verify}   & the reward model's own output (per-step) &
  training objective; labels (human, per-step) \\
\citet{laban2022summac}      & evaluation metric only & -- \\
\citet{song2024finesure}     & evaluation metric only & -- \\
\citet{min2023factscore}     & evaluation metric only & -- \\
\citet{zha2023alignscore}    & internal to the evaluator & -- \\
\midrule
This thesis & \textbf{construction of the preference label} &
  reward-model architecture and training procedure \\
\bottomrule
\end{tabularx}
\end{table}

Table~\ref{tab:rw-gap} organizes this literature by \emph{where} sentence-level
information is injected, which is the dimension along which the present study is
distinguished. Sentence-level factuality signals are clearly not unexplored: they
have been used to mask the policy-optimization objective
\citep{gu2025maskdpo}, to restructure the reward model
\citep{qiu2025sentencelevel}, to supply dense rewards during reinforcement
learning \citep{wu2023finegrained}, to train a reward model that outputs
per-step judgments directly \citep{lightman2023verify}, and to build
evaluation metrics \citep{laban2022summac, song2024finesure, min2023factscore}.
What distinguishes this thesis is the
insertion point. Each of the training-oriented works above takes the preference
labels as given, in most cases as response-level human judgments, and modifies
what happens downstream of them. The step examined here is the construction of
those labels: how the output of an automatic factuality evaluator, queried at
different granularities, is converted into a pairwise preference, and how
learnable the resulting labels are when everything downstream is held constant.
This is attractive because automatic labels are cheap, and uncertain because they
are noisier per label than human ones, so whether the additional granularity
recovers more signal than noise is not obvious in advance.

\section{Positioning of This Thesis}
\label{sec:rw-positioning}

The contribution of this thesis is not sentence-level reward modeling itself, nor
sentence-level factuality alignment in general, both of which are active areas
with established results. It is a controlled comparison of
\emph{preference-construction procedures}: the downstream Bradley--Terry
architecture and training procedure are held fixed while the granularity at which
automatic factuality evidence is converted into pairwise preference labels is
varied. No new factuality scorer is introduced; AlignScore is used off the shelf,
and the measurement instrument is the learnability of the resulting labels.

Six adjacent lines of work are distinguished as follows:

\begin{itemize}
  \item \textbf{RLAIF and LLM-as-a-judge research}
    \citep{bai2022constitutional, lee2023rlaif, zheng2023judging, gu2024survey}
    studies judge reliability at the level of a whole-response verdict, rather
    than the effect of querying the judge at a finer granularity.
  \item \textbf{Fine-grained human feedback} \citep{wu2023finegrained}
    establishes that localized supervision helps when labels are reliable; the
    present study asks the cheaper and less certain version, in which each local
    label is produced automatically.
  \item \textbf{Sentence-level factuality alignment} \citep{gu2025maskdpo}
    applies granular factuality information inside the policy-optimization
    objective while leaving preference construction unchanged; here the objective
    is unchanged and preference construction varies.
  \item \textbf{Sentence-level reward modeling} \citep{qiu2025sentencelevel}
    changes the reward model's architecture while learning from conventional
    response-level preferences; here the architecture is fixed and the
    preferences vary.
  \item \textbf{Process reward modeling} \citep{lightman2023verify} trains a
    reward model to output human-labeled, per-step correctness judgments
    directly, consumed downstream for reranking or search; here the reward
    model outputs a single scalar per summary, unchanged from the holistic
    condition, and what varies is only how that scalar's training label was
    constructed upstream.
  \item \textbf{Fine-grained evaluation frameworks}
    \citep{laban2022summac, song2024finesure, min2023factscore} decompose
    summaries in order to diagnose them rather than to supervise with them.
\end{itemize}

The comparison is conducted on preference data generated by one evaluator on one
corpus, so the findings describe behavior under specific conditions rather than a
general property of granular supervision
(Chapter~\ref{ch:discussion}, Section~\ref{sec:disc-generalizability}).

\chapter{Theoretical Background}
\label{ch:background}

\section{Preference Optimization: The Bradley--Terry Model}
\label{sec:bg-bradley-terry}

The Bradley--Terry model \citep{bradley1952rank} is a standard probabilistic model
of pairwise comparison.
It assumes that every item $i$ has an underlying, unobserved scalar quality $r_i$,
and that the probability of item $i$ being preferred over item $j$ in a pairwise
comparison is
\begin{equation}
  P(i \succ j) = \sigma(r_i - r_j),
  \qquad
  \sigma(x) = \frac{1}{1 + e^{-x}}.
  \label{eq:bradley-terry}
\end{equation}
Given a dataset of observed pairwise preferences $(w, l)$, where $w$ (the
``winner'') was judged preferable to $l$ (the ``loser''), the model parameters are
fit by maximizing the likelihood of the observed comparisons under
Equation~\ref{eq:bradley-terry}. Here the parameters belong to a neural network
$R_\theta$ that maps an (article, summary) pair to a scalar reward, and maximizing
the likelihood is equivalent to minimizing the negative log-likelihood
\begin{equation}
  \mathcal{L}(\theta) = - \, \mathbb{E}_{(w, l) \sim \mathcal{D}}
  \Big[ \log \sigma\big( R_\theta(w) - R_\theta(l) \big) \Big],
  \label{eq:bt-loss}
\end{equation}
where $\mathcal{D}$ is the preference dataset, holistic or GCA in this thesis.
Equation~\ref{eq:bt-loss} is exactly the loss implemented in \url{src/reward_model/train.py}
(\url{bradley_terry_loss}), and it belongs to the same family of
reward-modeling loss used inside RLHF pipelines more broadly
\citep{stiennon2020learning, ouyang2022training}. The difference here is that the
reward model $R_\theta$ is not an intermediate component on the way to fine-tuning a
generation policy; it is the final object of comparison. The research question this
thesis asks, whether GCA produces a more learnable preference signal than holistic
scoring, is operationalized directly: does a Bradley--Terry reward model trained on
GCA-derived preferences achieve higher held-out pairwise accuracy than one trained
on holistic-derived preferences, under an otherwise identical training recipe?

This corresponds to the inverse reinforcement learning perspective: given a set
of preference demonstrations, the task is to recover the latent reward function
that best explains them. Comparing two reward models trained on two differently
constructed preference datasets is then a controlled test of which construction
procedure yields the more consistently recoverable latent signal, without
requiring a downstream generation policy to be fine-tuned and re-evaluated
(Section~\ref{sec:bg-dpo}).

The architecture used throughout this thesis (\texttt{BradleyTerryRewardModel} in
\url{src/reward_model/train.py}) encodes the concatenation of a truncated source
article and a candidate summary with a pretrained transformer encoder, mean-pools
the encoder's final hidden states over non-padding tokens, and maps the pooled
representation to a scalar reward through a single linear layer. Reward-model
accuracy is defined as the fraction of held-out pairs for which the model correctly
ranks the chosen summary above the rejected one, i.e.\ $R_\theta(w) > R_\theta(l)$
(\texttt{\_evaluate} in \url{src/reward_model/train.py}); this is the metric
reported throughout Chapter~\ref{ch:results}.

\section{Direct Preference Optimization (DPO)}
\label{sec:bg-dpo}

Direct Preference Optimization (DPO) replaces RLHF's reward-model-plus-reinforcement-learning
pipeline with a single closed-form supervised objective on a policy, defined
relative to a fixed reference policy \citep{rafailov2023dpo}. It is included
here only as background on the study's earlier design: DPO was the objective
originally proposed for this thesis and produced the early results reported in
Appendix~\ref{app:exploratory}, Section~\ref{sec:results-early}, but it is not
used in the final experiments (Chapter~\ref{ch:methodology},
Section~\ref{sec:methodology-evolution}).

The relevant point is structural. The DPO objective consumes exactly the same
pairwise preference pairs used to train the reward models in
Section~\ref{sec:bg-bradley-terry}, but converts them directly into a policy
update, so the quality of the preference data and the behavior of the optimizer
both influence the resulting summaries. Separating those two influences is what
motivated measuring the preference data on its own, which is what the final
study does.

\section{AlignScore: Factual Consistency Scoring}
\label{sec:bg-alignscore}

AlignScore is a unified alignment function trained across a mixture of natural
language processing tasks recast into a single ``information alignment'' format
\citep{zha2023alignscore}. Given a context $c$, the source article in this thesis,
and a claim $a$, a candidate summary or a single sentence from one, AlignScore
produces a scalar score in $[0, 1]$ estimating the degree to which $a$ is supported
by $c$. This thesis uses AlignScore unmodified as the sole factuality judge for
both preference-construction regimes
(\url{src/judging/reward_model_judge.py}, class \texttt{RewardModelJudge}).

AlignScore exposes several evaluation modes, of which four are used in this thesis
(the \texttt{alignscore-evaluation-mode} argument of
\url{build_reward_preferences.py}):

\begin{itemize}
  \item \texttt{nli\_sp}: the context is first split into sentences before scoring,
    and the claim is scored against this sentence-split representation via
    AlignScore's NLI (three-way entailment/neutral/contradiction) output head. This
    was the default mode inherited from the earliest AlignScore-based experiments
    in the project.
  \item \texttt{nli}: the same NLI output head, without the sentence-splitting
    preprocessing step.
  \item \texttt{bin\_sp} / \texttt{bin}: the same two preprocessing variants, but
    using AlignScore's binary (aligned / not aligned) output head instead of the
    three-way NLI head.
\end{itemize}

Chapter~\ref{ch:results}, Section~\ref{sec:results-mode-sweep} reports an empirical
sweep across all four modes. The choice of mode is associated with substantial
changes in the outcome, alongside the GCA aggregation formula
(Chapter~\ref{ch:methodology}, Section~\ref{sec:bg-gca}). The \texttt{nli} mode produced the largest GCA
advantage in that sweep and was selected as the intended final
configuration for the validated $n=1{,}000$ campaign and the
scale-robustness checks (Appendix~\ref{app:exploratory},
Section~\ref{sec:results-final-campaign}), though
Chapter~\ref{ch:implementation}, Section~\ref{sec:impl-preferences} reports
that the confirmatory twenty-run campaigns were, due to an implementation
oversight, actually trained under \texttt{nli\_sp} instead. This should be read as a
judge-\emph{configuration} finding specific to this pipeline and this task, not as
a general property of sentence-level aggregation independent of the underlying
judge; Chapter~\ref{ch:discussion} discusses what can and cannot be concluded from
it.

\section{Statistical Validation: Bootstrap Confidence Intervals and the Wilcoxon Signed-Rank Test}
\label{sec:bg-stats}

Two statistical tools are used throughout Chapter~\ref{ch:results} to assess
whether an observed GCA-versus-holistic accuracy gap is distinguishable from
noise.

\textbf{Bootstrap confidence intervals.} The bootstrap \citep{efron1979bootstrap}
estimates the sampling distribution of a statistic by resampling the observed data.
Given a set of paired fold-level accuracy
differences $\{d_1, \ldots, d_m\}$, where $d_k = \mathrm{acc}_{\text{GCA}, k} -
\mathrm{acc}_{\text{Holistic}, k}$ for cross-validation fold $k$, pooled across
independent training seeds, a bootstrap 95\% confidence interval for the mean
difference $\bar{d}$ is constructed by resampling $\{d_1, \ldots, d_m\}$ with
replacement $B$ times, computing the mean of each resample, and taking the 2.5th
and 97.5th percentiles of the resulting distribution of means. This thesis uses $B
= 10{,}000$ resamples with a fixed NumPy random seed of 42.

\textbf{Wilcoxon signed-rank test.} As a complementary non-parametric test suited
to a small number of paired, non-normally-distributed differences, the two-sided
Wilcoxon signed-rank test \citep{wilcoxon1945individual} is reported on the same paired differences. Unlike a
paired $t$-test it does not assume normality, which is appropriate given the small
number of observations and the fact that fold accuracies are bounded proportions.

\textbf{An important qualification: the unit of replication.} Both procedures
assume that the observations they are applied to are independent, and this
determines which level the tests may legitimately be applied at. Fold-level
differences are \emph{not} independent: the five folds within a run partition a
single dataset and therefore share most of their training data, so applying
either procedure to 30 such observations treats correlated measurements as
independent replications, which understates uncertainty and produces intervals
that are too narrow and $p$-values that are too small. Run-level differences,
by contrast, come from separately seeded training runs and satisfy the
assumption.

This thesis therefore reports two kinds of figure, and the distinction should be
kept in view when reading Chapter~\ref{ch:results}. Fold-level intervals and
$p$-values, which appear only for the superseded six-run campaign
(Appendix~\ref{app:exploratory}, Section~\ref{sec:results-final-campaign}), are
anti-conservative and are presented as exploratory summaries of spread, never as
hypothesis tests. All confirmatory statistics (the twenty-run campaigns at each
dataset size, their bootstrap intervals, Wilcoxon tests, effect sizes, and
equivalence tests) are computed over run-level differences, where the
independence assumption holds. Chapter~\ref{ch:discussion},
Section~\ref{sec:disc-stats-caveat} sets out what each level can support.

\textbf{Effect size.} A $p$-value indicates whether a difference is
distinguishable from zero, not how large it is, so both are reported. For paired
differences $\{d_1, \ldots, d_n\}$ the standardized effect size is Cohen's $d =
\bar{d} / s_d$, the mean difference expressed in units of its own standard
deviation \citep{cohen1988statistical}. Alongside it, the matched-pairs
rank-biserial correlation is reported as the effect size native to the Wilcoxon
test: it is the difference between the sums of positive and negative signed
ranks, divided by their total, and therefore ranges from $-1$ (every run favors
one condition) to $+1$ (every run favors the other), capturing consistency of
direction independently of magnitude.

\textbf{Equivalence testing.} A non-significant result does not by itself
establish that an effect is absent; it may equally reflect insufficient data.
Establishing absence requires reversing the burden of proof, which is what the
two one-sided tests (TOST) procedure does \citep{lakens2017equivalence}. Given an
equivalence bound $\pm\delta$ representing the smallest effect that would be
practically meaningful, TOST tests the two null hypotheses that the true effect
is at least $+\delta$ and at most $-\delta$. Rejecting both places the true
effect inside $(-\delta, +\delta)$ with the stated confidence, which is positive
evidence of a negligible effect rather than mere failure to detect one. This
procedure is applied in Chapter~\ref{ch:results},
Section~\ref{sec:results-effect-sizes} to the $n=5{,}000$ and $n=10{,}000$
campaigns, where the substantive claim is precisely that no advantage remains.

Pooling across runs rather than reporting a single run is nonetheless deliberate,
and responds to an early observation (Appendix~\ref{app:exploratory},
Section~\ref{sec:results-alpha0-validation}): a promising single-run result, a
$+1.4$ percentage-point advantage from one hyperparameter configuration, did not
reproduce when the same configuration was run again. Single runs at this dataset
scale are therefore not reliable evidence on their own, whatever statistics are
computed from their folds.

\chapter{Methodology}
\label{ch:methodology}

\section{Research Questions and Hypotheses}
\label{sec:methodology-rqs}

\begin{itemize}
  \item \textbf{RQ1.} Does sentence-level factuality feedback aggregated via GCA
    produce a more learnable pairwise preference signal than holistic
    summary-level feedback, as measured by reward-model pairwise validation
    accuracy?

  \item \textbf{RQ2.} How do evaluator configuration and sentence-score
    aggregation strategy affect the relative learnability of holistic and
    GCA-derived preferences?

  \item \textbf{RQ3.} How stable are the observed differences across repeated
    reward-model training runs and dataset sizes?

  \item \textbf{RQ4.} Does the learnability advantage measured in RQ1 also
    hold when the trained reward models are scored against an independent,
    held-out ground truth constructed from the same factuality judge, rather
    than against the models' own training-style labels, and does that hold as
    the training-set size grows?
\end{itemize}

RQ1--RQ3 are addressed by the experiments in Chapter~\ref{ch:results} and
answered in Chapter~\ref{ch:conclusion}. RQ2 is addressed by an exploratory
sweep rather than a controlled ablation, so it supports statements about
association between configuration and outcome, not about causal mechanism
(Section~\ref{sec:methodology-exploratory}). RQ4 was added after the
$n=1{,}000$/$5{,}000$/$10{,}000$ campaign, in direct response to a limitation
that campaign left open: pairwise validation accuracy is measured against
labels supplied by the same evaluator that constructed the preferences, so a
reward model could in principle learn to recover those labels' idiosyncrasies
rather than genuine factuality signal (Section~\ref{sec:methodology-scope}).
Section~\ref{sec:methodology-ground-truth} describes the independent
evaluation constructed to address this, and
Section~\ref{sec:results-ground-truth} reports its results.

Three working hypotheses guided the design. H1: sentence-level aggregated
preferences are more learnable than holistic preferences, because they localize
the factuality signal within a long output. H2: the evaluator's configuration
affects whether any such difference is observable. H3: effects observed at one
dataset scale need not persist at another. Their status is assessed in
Chapter~\ref{ch:conclusion}.

\section{Study Design}
\label{sec:methodology-design}

The comparison is between two procedures for converting factuality scores into
pairwise preference labels. For every source article, two candidate summaries are
generated once and reused by both conditions:

\begin{description}
  \item[Holistic.] The complete candidate summary is scored against the source
    article, producing one score per candidate. The two scores are compared to
    yield a pairwise preference.
  \item[GCA.] The candidate summary is segmented into sentences; each sentence is
    scored separately against the source article; the sentence scores are
    aggregated into one summary-level score. The two aggregated scores are
    compared to yield a pairwise preference.
\end{description}

The following are held identical across conditions: the source articles; the two
candidate summaries and the decoding settings that produced them; the evaluator
family and checkpoint; the reward-model architecture; and the reward-model
training and evaluation procedure. The manipulated variable is the granularity at
which the evaluator is queried, together with the aggregation step that GCA
requires and holistic scoring does not.
Table~\ref{tab:methodology-summary} lists the final configuration.

\begin{table}[htbp]
\centering
\caption[Final experimental configuration]{Final experimental configuration.}
\label{tab:methodology-summary}
\begin{tabularx}{\textwidth}{lX}
\toprule
Factor & Value \\
\midrule
Task & Single-document abstractive summarization (news) \\
Dataset & CNN/DailyMail; runs at $n=1{,}000$, $5{,}000$, and $10{,}000$
  (subset seed 200) \\
Candidate generation & Mistral-7B-Instruct-v0.3, two-temperature sampling
  ($T=0.7$ / $T=1.0$) \\
Evaluator & \texttt{yzha/AlignScore}, mode \texttt{nli} \\
Manipulated variable & Feedback granularity: holistic vs.\ sentence-level $+$ GCA \\
GCA aggregation & Simple mean ($\alpha = 0.0$) \\
Margin & 0 (all pairs used) \\
Reward model & Bradley--Terry, \url{FacebookAI/roberta-base} backbone,
  mean-pool $+$ linear scalar head \\
RM training & epochs = 5, lr = $2\times10^{-5}$, batch size = 8,
  max length = 512, 5-fold cross-validation \\
Dependent variable & Reward-model pairwise validation accuracy \\
\bottomrule
\end{tabularx}
\end{table}

\section{Scope of the Evaluation}
\label{sec:methodology-scope}

The dependent variable throughout this thesis is the pairwise validation accuracy
of a Bradley--Terry reward model trained on a given preference set. Within this
thesis, \emph{learnability} denotes exactly this quantity: the extent to which the
preference relation encoded in a dataset can be recovered by the chosen
reward-model architecture under a fixed training procedure.

Because this term is easily over-read, it is worth stating precisely what a higher
accuracy does and does not establish. A higher value indicates that the generated
preference relation is more consistently recoverable by the selected reward-model
architecture. It does not by itself establish any of the following:

\begin{itemize}
  \item that the preference labels agree better with human factuality judgments;
  \item that the labels are objectively more accurate as factuality annotations;
  \item that a summarization model trained on those preferences produces more
    factually consistent summaries;
  \item that GCA is generally preferable to holistic scoring.
\end{itemize}

The first two are inaccessible here because no independent human annotation was
collected, so the labels produced by AlignScore \citep{zha2023alignscore} function
as the reference standard rather than as verified ground truth. Substituting an
automatic judge for human annotators is the defining move of RLAIF
\citep{bai2022constitutional, lee2023rlaif}, and it carries this limitation with
it: the resulting labels are cheap and consistent, but unaudited. The third is inaccessible because the final
study does not include a downstream policy-optimization stage
(Section~\ref{sec:methodology-evolution}). The fourth is addressed empirically and
answered in the negative: Chapter~\ref{ch:results} reports configurations in which
GCA does not lead. Claims made in later chapters are confined to learnability in
the sense defined here.

\section{Reward-Model Evaluation Protocol}
\label{sec:methodology-rm-eval}

Reward-model accuracy is evaluated by 5-fold cross-validation
\citep{kohavi1995crossvalidation} rather than a single fixed held-out split. At
the dataset sizes used here, a fixed split would leave either too little data
for training or too little for a stable accuracy estimate, whereas
cross-validation uses the full dataset for both and yields five accuracy
readings per run.

Cross-validation alone does not account for variance introduced by the reward
model's initialization and by which pairs fall into which fold. This became
apparent during the study: a first run of the final configuration produced a GCA
advantage of $+6.0$ percentage points, while a rerun of the identical
configuration \emph{at the same seed value} produced $+1.3$ percentage points
(Chapter~\ref{ch:results}, Section~\ref{sec:results-mode-sweep}). That two
same-seed runs diverge follows from an implementation detail documented in
Chapter~\ref{ch:implementation}, Section~\ref{sec:impl-rm-training}: the
\texttt{--seed} argument controls the cross-validation fold assignment, while
PyTorch's generator is left unseeded, so weight initialization, batch ordering and
dropout all vary between executions. A seed value therefore does not pin down a
run completely, and repeated runs at one seed value are not redundant.

The $n=1{,}000$ result was accordingly first pooled over \emph{six training runs}
spanning \emph{five distinct seed values} (42, 42, 7, 100, 314, 2026; seed 42
appears twice, as the original run and its rerun), giving 30 fold-level accuracy
readings per condition.

Two properties of this pooling constrain the statistical claims that can be made
from it, and are revisited in Chapter~\ref{ch:discussion},
Section~\ref{sec:disc-stats-caveat}. Six runs across five seed values is a modest
basis for a variance estimate. More consequentially, the 30 fold-level differences
are not 30 independent observations: the five folds within a run partition a
single dataset, so their training sets overlap heavily and their errors are
correlated. Procedures assuming independent observations therefore overstate the
effective sample size. Fold-level uncertainty estimates are reported in
Chapter~\ref{ch:results} as exploratory and anti-conservative, not as
confirmatory tests.

Both constraints motivated the design actually used for the reported headline
result. The \emph{run} is treated as the unit of independent replication
throughout, and the campaign was extended to \emph{twenty independently seeded
runs} (seeds 1--20, fixed before any was executed) under a corrected training
procedure in which PyTorch and CUDA are seeded per fold, so that a run is fully
determined by its seed (Chapter~\ref{ch:implementation},
Section~\ref{sec:impl-rm-training}). All confirmatory statistics reported in
this thesis --- the bootstrap interval, the Wilcoxon test, the effect sizes, and
the equivalence tests of Section~\ref{sec:results-effect-sizes} --- are computed
over these run-level differences, for which the independence assumption does
hold. Fold-level figures are retained only as descriptive summaries of spread.
Twenty runs is also the smallest campaign size that permits a two-sided Wilcoxon
result below the conventional threshold with a realistic margin: with $n$ paired
same-signed observations the smallest attainable two-sided $p$-value is
$2/2^{n}$, which at $n=6$ is $0.031$ and leaves no room for a single
discordant run, while at $n=20$ it is $1.9\times10^{-6}$
(Chapter~\ref{ch:results}, Section~\ref{sec:results-extended-campaign}).

The $n=5{,}000$ and $n=10{,}000$ configurations were initially each executed
once rather than repeated, for reasons of computational budget, and at that
stage carried correspondingly less weight than the pooled $n=1{,}000$ result.
Both were subsequently extended to twenty-run campaigns under the same
corrected, deterministic training procedure used for the $n=1{,}000$ extension,
matching its resolution exactly (Chapter~\ref{ch:results},
Section~\ref{sec:results-scale}), so all three dataset sizes are now supported
by equally powered, repeated-run evidence rather than only $n=1{,}000$.

\section{Exploratory and Confirmatory Experiments}
\label{sec:methodology-exploratory}

The experiments in Chapter~\ref{ch:results} do not all have the same evidential
status, and the distinction is material to how the results should be read.

The aggregation-formula comparison (Section~\ref{sec:results-formula-opt}), the
evaluator-mode sweep (Section~\ref{sec:results-mode-sweep}), and the
hyperparameter search (Section~\ref{sec:results-alpha0-validation}) were all
carried out on the $n=1{,}000$ setting, and the final configuration was selected
after observing their outcomes. These experiments are exploratory. The
configuration they produced was not specified in advance, and reporting it as
though it had been pre-registered would overstate the evidence: with enough
configurations examined, some separation between conditions is expected by chance
alone.

The repeated runs at $n=1{,}000$ (Section~\ref{sec:results-final-campaign}) test
whether the selected configuration produces a stable effect under repetition, but
they reuse the dataset on which that configuration was chosen and therefore cannot
establish that it generalizes beyond that setting.

The $n=5{,}000$ and $n=10{,}000$ experiments provide a weaker check than would be
ideal, for a reason that must be stated explicitly. All three subsets are drawn
from the same CNN/DailyMail test split by the same deterministic procedure at the
same seed, differing only in the number of samples requested, and they are
consequently nested rather than disjoint
(Section~\ref{sec:methodology-nesting}). The larger-scale runs therefore
re-use almost all of the data on which the configuration was selected, and are
best understood as measurements of how the effect behaves as the dataset grows,
not as out-of-sample replication on independent data. No experiment in this thesis
evaluates the selected configuration on data disjoint from the set used to choose
it.

\section{Relationship Between the Dataset Sizes}
\label{sec:methodology-nesting}

The three dataset sizes used in the final experiments are not independent samples,
and this constrains how the larger-scale runs can be interpreted.

All three subsets are produced by the same deterministic procedure
(\url{src/data/subset.py}, \texttt{select\_subset}) from the same CNN/DailyMail
test split, under the same length filters and the same subset seed (200). Only the
number of requested samples differs. Because the selection draws from one eligible
pool of 10{,}896 articles with a fixed random seed, the resulting subsets overlap
heavily. Reproducing the selection with the recorded parameters gives the
intersections in Table~\ref{tab:subset-overlap}.\footnote{Computed by re-running
\texttt{select\_subset} with the parameters recorded in
\url{configs/subset_1000.yaml} and \url{slurm/generate_candidates_scale.sh}
($\texttt{seed}=200$, test split, article limit 8{,}000 characters, summary limit
2{,}000 characters) against the local copy of the test split in
\url{data/cnn_dailymail_test_full}, and comparing the resulting SHA-256-derived
sample IDs. The candidate and preference files for these runs are not retained in
the repository, so the comparison reproduces the selection rather than reading the
stored artifacts; the selection is deterministic given these parameters.}

\begin{table}[htbp]
\centering
\caption[Overlap between the three dataset subsets]{Overlap between the subsets used in the final experiments. All three use
subset seed 200 and differ only in sample count.}
\label{tab:subset-overlap}
\begin{tabular}{lrr}
\toprule
Pair & Shared sample IDs & Share of the smaller set \\
\midrule
$n=1{,}000$ and $n=5{,}000$   & 980   & 98.0\% \\
$n=1{,}000$ and $n=10{,}000$  & 999   & 99.9\% \\
$n=5{,}000$ and $n=10{,}000$  & 5{,}000 & 100\% \\
\bottomrule
\end{tabular}
\end{table}

\begin{figure}[htbp]
\centering
\begin{tikzpicture}[font=\footnotesize]
  % eligible pool
  \draw[gray!55, fill=gray!8, rounded corners=2pt] (0,0) rectangle (11.2,4.3);
  \node[anchor=north west, gray!55!black, font=\scriptsize]
        at (0.2,4.15) {eligible pool: 10{,}896 articles};

  % n = 10000
  \draw[draw=orange!65!black, fill=orange!10, rounded corners=2pt]
        (0.45,0.35) rectangle (10.75,3.55);
  \node[anchor=north west, orange!55!black, font=\scriptsize]
        at (0.62,3.42) {$n=10{,}000$};

  % n = 5000
  \draw[draw=blue!60!black, fill=blue!10, rounded corners=2pt]
        (0.85,0.7) rectangle (5.95,2.95);
  \node[anchor=north west, blue!50!black, font=\scriptsize]
        at (1.02,2.82) {$n=5{,}000$};

  % n = 1000
  \draw[draw=black!70, fill=white, rounded corners=2pt]
        (1.25,1.05) rectangle (2.85,2.25);
  \node[anchor=north west, font=\scriptsize] at (1.4,2.12) {$n=1{,}000$};

  % annotations with leader lines, placed in free space
  \node[anchor=west, font=\scriptsize, black!75] (a1) at (6.35,2.55)
       {980 of 1{,}000 also in $n{=}5{,}000$};
  \draw[black!45, -{Stealth[length=1.5mm]}] (a1.west) -- (2.95,2.0);

  \node[anchor=west, font=\scriptsize, black!75] (a2) at (6.35,1.85)
       {999 of 1{,}000 also in $n{=}10{,}000$};
  \draw[black!45, -{Stealth[length=1.5mm]}] (a2.west) -- (2.95,1.55);

  \node[anchor=west, font=\scriptsize, blue!50!black] (a3) at (6.35,1.15)
       {all 5{,}000 also in $n{=}10{,}000$};
  \draw[blue!40, -{Stealth[length=1.5mm]}] (a3.west) -- (6.0,1.0);
\end{tikzpicture}
\caption[Nesting of the three dataset subsets]{The three subsets are nested rather than independent. Box sizes are
schematic; the counts are the measured intersections reported in
Table~\ref{tab:subset-overlap}.}
\label{fig:subset-nesting}
\end{figure}

The $n=5{,}000$ subset is contained entirely within the $n=10{,}000$ subset, and
the $n=1{,}000$ subset is almost entirely contained within both. The
$n=10{,}000$ set additionally covers 91.8\% of the whole eligible pool, so at that
size little of the test split remains unused in any case.

Two consequences follow, and both are carried through the rest of the thesis. The
larger-scale experiments do not constitute out-of-sample replication: they re-use
almost all of the data on which the configuration was selected, and therefore
measure how the effect behaves as further data is added to the same pool rather
than whether it transfers to new data. And because the samples are nested, the
three results are not statistically independent of one another, so the differences
between them should not be read as three separate estimates of the same quantity.
Drawing the larger subsets with different seeds would have made them independent,
at the cost of the nesting property the generation script was written to provide;
this is recorded as a limitation in Chapter~\ref{ch:conclusion}.

\section{Independent Ground-Truth Construction}
\label{sec:methodology-ground-truth}

The evaluation described above has two properties that limit what it can
establish, both noted in Section~\ref{sec:methodology-exploratory} and
Section~\ref{sec:concl-limitations}: the labels a trained reward model is
scored against are supplied by the same evaluator that constructed its
training preferences, and the $n=5{,}000$/$n=10{,}000$ checks reuse almost
all of the data the configuration was selected on rather than evaluating it
on disjoint data. RQ4 was formulated to address both at once, by scoring the
already-trained reward models against a held-out set built specifically to
be disjoint from every training pool used in this thesis, under an
evaluation criterion that does not depend on how any individual training
pair happened to be labelled.

\paragraph{Held-out article selection.} 500 articles were drawn from the same
CNN/DailyMail test split, using \texttt{select\_disjoint\_subset}
(\url{src/data/subset.py}): the function reconstructs the exact
\texttt{rng.sample()} call that produced the $n=10{,}000$, seed-200 draw used
throughout this thesis, removes those indices from the eligible pool, and
draws 500 samples from what remains under a new seed (999). Because the
$n=5{,}000$ subset is fully contained in the $n=10{,}000$ draw and the
$n=1{,}000$ subset is 99.9\% contained in it
(Table~\ref{tab:subset-overlap}), excluding the $n=10{,}000$ draw is
sufficient to make the held-out set disjoint from all three training pools
used in the main campaign; two further training pools introduced later in
this work, at $n=2{,}000$ and $n=3{,}000$ (Section~\ref{sec:results-scale-curve}),
were verified by direct index comparison to be exact subsets of the
$n=10{,}000$ draw and are therefore disjoint from the held-out set for the
same reason.

\paragraph{Candidate generation and scoring.} Two candidate summaries were
generated per held-out article with the same Mistral-7B-Instruct-v0.3 pipeline
and the same two decoding temperatures used throughout this thesis
(Section~\ref{sec:methodology-design}), giving 1{,}000 candidate summaries in
total. Every candidate was scored with AlignScore under the locked
configuration used for the main campaign: sentence-level scoring, GCA
aggregation with $\alpha=0.0$ (Section~\ref{sec:results-formula-opt}). This
GCA-aggregated score is used only to construct the evaluation criterion
below; it is never used as a training signal for the held-out set, and no
reward model is trained on these 500 articles.

\paragraph{Evaluation criterion.} Two ground-truth constructions were used.
The first pairs each article's own two candidates directly, mirroring the
structure of the training preferences; the two candidates are generated by
the same model at similar temperatures, so their quality gap is typically
small (mean GCA score gap approximately 0.04). The second pools all
1{,}000 scored candidates across all 500 articles, sorts them by GCA score,
and compares the top $X\%$ against the bottom $X\%$ for $X\in\{25,10,5\}$,
which produces progressively larger, more unambiguous quality gaps between
the two sides at the cost of fewer summaries per comparison
(Section~\ref{sec:results-ground-truth}). At $X=5$, every high-scoring
summary is compared against every low-scoring summary (2{,}500 pairs from 50
summaries per side) rather than one random pairing, to remove pairing-order
noise from the smallest, highest-signal condition. A trained reward model's
\emph{success rate} on a given construction is the fraction of comparisons in
which it scores the higher-GCA-score summary above the lower-GCA-score one;
this is the dependent variable throughout Section~\ref{sec:results-ground-truth},
and it is defined identically for the holistic and GCA reward models, so
neither construction favours either condition by design.

\paragraph{Checkpoint retention.} The training procedure used for RQ1--RQ3
(Section~\ref{sec:methodology-rm-eval}) discards each fold's trained weights
after recording its accuracy, since only the accuracy is needed there. RQ4
requires scoring the same trained model against a second, independent
criterion, so the training procedure was changed for this evaluation to save
the full reward-model checkpoint (encoder weights and reward head) for every
run, rather than only its cross-validation accuracy.

\section{Evolution of the Study Design}
\label{sec:methodology-evolution}

The study reported here is narrower in scope than the one outlined in the original
research plan, and the reasons are methodological.

The initial design used the generated preferences directly for DPO-based policy
optimization, evaluating factual consistency on the fine-tuned model's own
generations. Early downstream experiments made it difficult to attribute observed
differences to any single stage: a change in generated-summary factuality could
originate in the preference-construction procedure, in the quality of the reward
signal, or in the policy-optimization step, and the design provided no way to
separate these. Increasing the number of stages also increased the number of
variance sources standing between the manipulated variable and the measured
outcome.

The final study therefore isolates an earlier and more controlled diagnostic
question: whether the preference data produced by each construction procedure is
recoverable by a fixed reward-model architecture. Removing the policy-optimization
stage removes both reinforcement-learning dynamics and decoding variance from the
measurement path, leaving a comparison in which granularity is the principal
difference between conditions.

The claims of this thesis are correspondingly narrower. The work evaluates the
learnability of the preference signal rather than demonstrating improvements in
generated-summary factuality (Section~\ref{sec:methodology-scope}). Two further
consequences follow from the same change. Margin-based filtering of near-tie pairs
was removed, since AlignScore produces continuous scores for which exact ties are
effectively impossible, and the threshold discarded roughly one fifth of the data
without a demonstrated benefit. The reliability protocol designed for a generative
LLM judge, comprising confidence gating, order randomization, and rationale
requirements, no longer applies once the evaluator is a deterministic model that
is never prompted and returns no rationale; a planned human audit of judge
rationales was not carried out and is recorded as a limitation in
Chapter~\ref{ch:conclusion}.

\section{Validity Considerations}
\label{sec:methodology-ethics}

\paragraph{Ethical and practical considerations.} All experiments use
CNN/DailyMail \citep{hermann2015teaching, nallapati2016abstractive}, a publicly
available corpus of news articles and human-written summaries; no personally
sensitive data is collected. Because both the evaluator's
judgments and the generated summaries can be incorrect, results are reported as
measurements of preference-data learnability rather than as statements about
factual correctness, and model outputs are not presented as verified fact. Compute
use was bounded by working with fixed-size subsets.

\paragraph{Candidate generation and decoding temperature.} The two candidates in
each pair are produced by the same model at different decoding temperatures
($T=0.7$ and $T=1.0$). This introduces a potential confound: preference structure
may correlate with systematic stylistic differences induced by temperature rather
than with factuality alone. The preference data stored in the repository allow the
asymmetry to be quantified directly, and Table~\ref{tab:temperature-asymmetry}
reports it.\footnote{Computed from
\url{data/preferences/holistic_reward_preferences_200.jsonl} and
\url{data/preferences/gca_reward_preferences_200.jsonl}. The margin-0 rows
recompute the decision from the stored scores at the threshold used in the final
configuration; the margin-0.05 rows reproduce the stored decisions, whose
aggregate counts also appear in
\url{reports/reward_model_judging_summary.csv}. Preference files for the final
$n=1{,}000$ configuration are not retained in the repository, so these proportions
could not be recomputed under the final evaluator mode.}

\begin{table}[htbp]
\centering
\caption[Decoding-temperature asymmetry in the preference labels]{Proportion of pairs in which the lower-temperature candidate ($T=0.7$)
is preferred, on the $n=200$ preference set (evaluator mode \texttt{nli\_sp},
$\alpha=0.5$).}
\label{tab:temperature-asymmetry}
\begin{tabular}{llrrr}
\toprule
Condition & Threshold & $T{=}0.7$ preferred & $T{=}1.0$ preferred & Share \\
\midrule
Holistic & margin $=0$    & 137 & 63 & 68.5\% \\
GCA      & margin $=0$    & 126 & 74 & 63.0\% \\
\addlinespace
Holistic & margin $=0.05$ & 109 & 45 & 70.8\% \\
GCA      & margin $=0.05$ &  97 & 60 & 61.8\% \\
\bottomrule
\end{tabular}
\end{table}

The asymmetry is substantial under both conditions and consistently larger under
holistic scoring, and it is not an artifact of the margin threshold. Two
implications follow. Features correlated with decoding temperature, such as
summary length or lexical conservativeness, may carry a non-trivial share of the
label information, so a reward model could recover part of the preference relation
without representing factuality as such. And because the asymmetry differs between
conditions by roughly five percentage points, the two preference sets are not
equivalent in this respect either, which means the comparison between them is not
perfectly clean with respect to temperature.

These figures come from the $n=200$ set scored under the earlier evaluator
configuration and need not hold exactly at the final one. They are reported as
evidence that the confound is real and measurable, not as a quantification of its
effect on the reported results. Controlling candidate generation more explicitly,
by sampling both candidates at one temperature, would remove the confound at
source and is left to future work (Chapter~\ref{ch:conclusion}).

\paragraph{How much label information do surface features actually carry?} The
concern above is only as serious as the amount of label information that
temperature-correlated surface features actually carry, and that is measurable
directly on the same stored data. Table~\ref{tab:surface-baselines} reports the
accuracy of the cheapest possible classifiers --- ones that ignore the article
entirely and decide from a single surface property of the two candidate
summaries --- against the same decisions the evaluator
made.\footnote{Computed by \url{analysis/surface_feature_baselines.py} from the
same two preference files as Table~\ref{tab:temperature-asymmetry}, restricted
to the pairs carrying an A/B decision (154 holistic, 157 GCA).}

\begin{table}[htbp]
\centering
\caption[Surface-feature baselines against the judge's decisions]{Accuracy of surface-feature classifiers against the evaluator's own
decisions, $n=200$ preference set (mode \texttt{nli\_sp}, $\alpha=0.5$). Each
row states the accuracy of the better of the two opposing rules (prefer
longer/shorter, prefer more/fewer sentences). Mean lengths are in characters.}
\label{tab:surface-baselines}
\begin{tabular}{lrr}
\toprule
Predictor & Holistic & GCA \\
\midrule
Prefer the longer summary        & 0.494 & 0.465 \\
Prefer the shorter summary       & 0.506 & 0.535 \\
Prefer more sentences            & 0.519 & 0.475 \\
Prefer fewer sentences           & 0.481 & 0.525 \\
\addlinespace
Mean length, chosen              & 865.6 & 840.5 \\
Mean length, rejected            & 863.5 & 846.2 \\
\bottomrule
\end{tabular}
\end{table}

None of these predictors performs meaningfully above chance. The mean length of
the chosen summary differs from that of the rejected summary by 2 characters
under holistic scoring and by $-6$ characters under GCA, on summaries averaging
roughly 850 characters, and sentence counts differ by 0.01 and $-0.28$
respectively. Length and sentence count, the two most obvious surface proxies
for decoding temperature, therefore do not carry the preference signal, and a
reward model cannot be recovering its accuracy from them.

This narrows the confound without eliminating it, and the distinction matters.
The 68.5\% and 63.0\% asymmetries in
Table~\ref{tab:temperature-asymmetry} are not themselves a shortcut available to
the reward model: the model is trained on (article, chosen) and (article,
rejected) text pairs and never observes which candidate came from which
temperature, or which position it occupied
(Chapter~\ref{ch:implementation}, Section~\ref{sec:impl-rm-training}). For the
asymmetry to inflate reward-model accuracy, some \emph{textual} property
correlated with temperature would have to be learnable from the summary itself.
Table~\ref{tab:surface-baselines} rules out the two most likely candidates for
such a property. Subtler lexical or stylistic correlates are not ruled out, and
testing them would require the feature-level analysis recorded as future work,
but the simple version of the objection does not hold.

\paragraph{Threats to validity.} Table~\ref{tab:threats-to-validity} summarizes
the threats considered and how each is handled.

\begin{table}[htbp]
\centering
\caption[Threats to validity]{Threats to validity and how each is addressed.}
\label{tab:threats-to-validity}
\begin{tabularx}{\textwidth}{lX}
\toprule
Threat & Status \\
\midrule
Single evaluator (AlignScore) & Not mitigated; all results are conditional on this
  evaluator's judgments, which are not independent ground truth. \\
Evaluator and aggregation sensitivity & Measured directly via the formula
  comparison and evaluator-mode sweep (Chapter~\ref{ch:results}). \\
Fold- and run-level noise & Addressed by twenty independently seeded runs at each
  of the three dataset sizes, under fully deterministic per-fold seeding
  (Section~\ref{sec:methodology-rm-eval}). \\
Non-independence of pooled folds & Mitigated by testing at the run level, where
  the independence assumption holds; fold-level figures are reported only as
  descriptive summaries of spread (Section~\ref{sec:methodology-rm-eval}). \\
Configuration selected on observed results & Acknowledged, not remedied; the
  larger-scale subsets are nested supersets of the selection data
  (Section~\ref{sec:methodology-nesting}), so no run evaluates the configuration
  on independent data. \\
Scale sensitivity & Examined explicitly rather than assumed away
  (Section~\ref{sec:results-scale}). \\
Decoding-temperature confound & Quantified where data permit; not controlled for
  in the design. Length and sentence count are excluded as carriers of the
  signal (Table~\ref{tab:surface-baselines}); subtler lexical correlates remain
  untested (see above). \\
Source-document truncation & Documented in Chapter~\ref{ch:implementation},
  Section~\ref{sec:impl-rm-training}; discussed as a limitation. \\
Weak absolute accuracy & Both conditions remain near chance
  (Chapter~\ref{ch:discussion}, Section~\ref{sec:disc-weak-signal}). \\
\bottomrule
\end{tabularx}
\end{table}

\chapter{Implementation}
\label{ch:implementation}

\section{Pipeline Overview}
\label{sec:impl-pipeline}

Figure~\ref{fig:pipeline} shows the final pipeline. It differs from the pipeline
described in the original proposal in exactly the ways described in
Chapter~\ref{ch:methodology}: there is no DPO fine-tuning stage, no MoDPO branch,
and no cross-summary sentence-alignment step. The pipeline has four stages. First,
a fixed number of CNN/DailyMail articles are selected deterministically
(Section~\ref{sec:impl-data}). Second, two candidate summaries are generated per
article at different sampling temperatures (Section~\ref{sec:impl-generation}).
Third, the same candidate pool is scored twice, once holistically and once at the
sentence level with GCA aggregation, to produce two separate preference datasets
(Section~\ref{sec:impl-preferences}). Fourth, a Bradley--Terry reward model is
trained independently on each preference dataset and the two are compared
(Section~\ref{sec:impl-rm-training}).

\begin{figure}[htbp]
\centering
\begin{tikzpicture}[
  node distance=7mm and 8mm,
  box/.style={draw, rounded corners, align=center, font=\footnotesize,
    minimum width=34mm, minimum height=9mm, inner sep=2mm},
  shared/.style={box, fill=gray!12},
  hol/.style={box, fill=blue!8},
  gca/.style={box, fill=orange!10},
  end/.style={box, fill=gray!12, minimum width=48mm},
  arr/.style={-{Stealth[length=2mm]}, thick},
]

% shared stages
\node[shared] (data)   {1. Data Loader\\\footnotesize CNN/DailyMail, seeded subset};
\node[shared, below=of data]  (gen)   {2. Candidate Generation\\\footnotesize Mistral-7B, $T{=}0.7$ / $T{=}1.0$};
\node[shared, below=of gen]   (pair)  {3. Candidate Pairing\\\footnotesize Summary A, Summary B};

% branch anchor row
\node[hol, below left=12mm and -8mm of pair]  (hscore) {4A. Holistic Scoring\\\footnotesize AlignScore on full summary};
\node[gca, below right=12mm and -8mm of pair] (seg)    {4B. Sentence Segmentation};

\node[hol, below=of hscore] (hpref)  {5A. Holistic Preference\\\footnotesize margin $=0$};
\node[gca, below=of seg]    (gscore) {5B. Sentence-Level Scoring\\\footnotesize AlignScore per sentence};

\node[gca, below=of gscore] (gagg)   {6B. GCA Aggregation\\\footnotesize $\alpha{=}0.0$ (simple mean)};
\node[gca, below=of gagg]   (gpref)  {7B. GCA Preference\\\footnotesize margin $=0$};

\node[hol, below=18mm of hpref] (hrm) {8A. Bradley--Terry RM\\\footnotesize \emph{RM-Holistic}, 5-fold CV};
\node[gca, below=of gpref]      (grm) {8B. Bradley--Terry RM\\\footnotesize \emph{RM-GCA}, 5-fold CV};

\node[end, below=12mm of grm] (eval) {9. Evaluation and Comparison\\\footnotesize pairwise accuracy, bootstrap CI, Wilcoxon test};

% arrows
\draw[arr] (data) -- (gen);
\draw[arr] (gen) -- (pair);
\draw[arr] (pair) -| (hscore);
\draw[arr] (pair) -| (seg);
\draw[arr] (hscore) -- (hpref);
\draw[arr] (seg) -- (gscore);
\draw[arr] (gscore) -- (gagg);
\draw[arr] (gagg) -- (gpref);
\draw[arr] (hpref) -- (hrm);
\draw[arr] (gpref) -- (grm);
\draw[arr] (hrm) |- (eval);
\draw[arr] (grm) -- (eval);

\end{tikzpicture}
\caption[The final experimental pipeline]{The final pipeline. Shared setup (grey) produces one candidate pool from
which two independent preference-construction branches run: holistic (blue) and
GCA (orange). Each branch trains its own Bradley--Terry reward model under an
identical recipe, and the two are compared in the final evaluation step. Compare
to the pipeline described in the original proposal (Chapter~\ref{ch:methodology}),
which additionally included a DPO fine-tuning stage and a cross-summary
sentence-alignment step, neither of which appears here.}
\label{fig:pipeline}
\end{figure}

\section{Data}
\label{sec:impl-data}

All experiments use CNN/DailyMail \citep{hermann2015teaching}, version \texttt{3.0.0}
as distributed through the Hugging Face \texttt{datasets} library. Subset selection
(\url{src/data/subset.py}, function \texttt{select\_subset}) filters the chosen
split to articles under 8{,}000 characters and reference summaries under 2{,}000
characters, then draws a fixed-size sample from the eligible pool using Python's
\texttt{random.Random} seeded deterministically. Each sample is assigned an ID of
the form \texttt{cnn\_\{index\}\_\{hash\}}, where \texttt{hash} is the first 12
hexadecimal characters of the SHA-256 digest of the article text
(\texttt{make\_sample\_id}). This ties every sample's identifier to its content
rather than to its position in the dataset, so the same article always receives the
same ID regardless of which subset size or seed selected it.

Table~\ref{tab:impl-subsets} lists every subset used across the project. The
$n=1{,}000$, $n=5{,}000$, and $n=10{,}000$ subsets all use subset seed 200 and the
same filtering parameters, differing only in the number of samples requested. They
are therefore \emph{nested rather than disjoint}: the smaller sets are almost
entirely contained in the larger ones. Chapter~\ref{ch:methodology},
Section~\ref{sec:methodology-nesting} quantifies the overlap and states its
methodological consequence. Earlier subsets at seeds 42 and 100 were drawn with
different seeds and are not part of the final comparison.

\begin{table}[htbp]
\centering
\caption[CNN/DailyMail subsets used across the project]{CNN/DailyMail subsets used across the project.}
\label{tab:impl-subsets}
\begin{tabularx}{\textwidth}{lclX}
\toprule
$n$ & Subset seed & Purpose & Notes \\
\midrule
200    & 42  & Early DPO-based prototype & Superseded; see
  Chapter~\ref{ch:results}, Section~\ref{sec:results-early} \\
495    & 100 & First Bradley--Terry RM training set & Drawn with a different seed
  from the $n=200$ set \\
1{,}000 & 200 & Final repeated-run campaign & 6 RM training runs; seeds 42, 42, 7, 100,
  314, 2026 \\
5{,}000 & 200 & Scale-robustness check & Single training run \\
10{,}000 & 200 & Scale-robustness check & Single training run \\
\bottomrule
\end{tabularx}
\end{table}

\section{Candidate Generation}
\label{sec:impl-generation}

Candidate summaries are generated with Mistral-7B-Instruct-v0.3, loaded in
\texttt{bfloat16} without quantization (\url{src/generation/candidates.py}). On the
A100-SXM4-40GB GPUs used for this project, the model occupies roughly 14\,GB in
\texttt{bfloat16}, comfortably within budget, so 4-bit quantization, though
supported in the code via \texttt{BitsAndBytesConfig}, is disabled in the
configuration actually used (\texttt{load\_in\_4bit: false} in
\url{configs/generation.yaml}). The tokenizer uses left-padding, required for
batched generation with a decoder-only model, and articles are truncated to 1{,}024
input tokens.

Each article is summarized twice, using the fixed prompt template
\texttt{"Summarize the following news article in a concise paragraph.\ \ldots"},
with two different sampling configurations: a low-temperature setting
($T=0.7$, top-$p=0.9$) and a high-temperature setting ($T=1.0$, top-$p=0.95$),
generating up to 256 new tokens each. These two outputs form the candidate pair
(``summary A'' and ``summary B'') that both the holistic and GCA preference
construction stages score. Generation is resumable: the script checks which
sample IDs already appear in the output file and only generates candidates for the
remainder, which matters in practice on a shared cluster where a job may need to be
resubmitted after a queue timeout.

\section{Preference Construction}
\label{sec:impl-preferences}

Both preference-construction regimes are implemented in
\url{src/judging/build_reward_preferences.py} and share a single judge object,
\texttt{RewardModelJudge} (\url{src/judging/reward_model_judge.py}), configured to
use the AlignScore backend.

In the holistic regime (\texttt{judge\_holistic\_pair}), each full candidate
summary is scored once against the article with \texttt{judge.score(article,\
summary)}, and the two resulting scores are compared with a margin guard
(\texttt{make\_decision}): summary A is preferred if its score exceeds summary B's
by more than the margin, summary B is preferred in the opposite case, and the pair
is marked \texttt{no\_preference} otherwise. With the final margin of 0
(Chapter~\ref{ch:methodology}, Section~\ref{sec:methodology-evolution}), every
pair with a nonzero score
difference produces a decision.

In the GCA regime (\texttt{judge\_gca\_pair}), each candidate summary is first
split into sentences with a regex-based splitter, then every sentence is scored
independently against the article with \texttt{judge.score\_sentences(article,\
sentences)}. The resulting list of per-sentence scores is passed to
\texttt{aggregate\_sentence\_scores(scores,\ alpha)}
(Chapter~\ref{ch:background}, Equation~\ref{eq:gca-general}) to produce a single
GCA score per summary, and the same margin-guarded comparison determines the
preference. Every output record retains the full sentence-level detail, including
each sentence's individual AlignScore score and a count of ``low-faithfulness''
sentences (score below 0.5), which supports the qualitative disagreement analysis
in Chapter~\ref{ch:results}.

The locked final configuration is invoked as:
\begin{quote}
\small\ttfamily
python src/judging/build\_reward\_preferences.py\\
\hspace*{1em}--mode both --alpha 0.0 --margin 0\\
\hspace*{1em}--alignscore-evaluation-mode nli\\
\hspace*{1em}--alignscore-backbone FacebookAI/roberta-base
\end{quote}

\section{Reward-Model Training}
\label{sec:impl-rm-training}

Reward-model training is implemented in \url{src/reward_model/train.py}
(\texttt{BradleyTerryRewardModel}, \texttt{bradley\_terry\_loss}) and orchestrated
by \url{src/reward_model/run_training.py}, which trains the holistic and GCA
conditions from a single invocation. All results reported in
Chapter~\ref{ch:results} use $k$-fold cross-validation
(\texttt{\_kfold\_cv} in \texttt{run\_training.py}): the preference pairs are
shuffled using a seeded \texttt{random.Random} instance, split into $k=5$
contiguous folds, and for each fold a freshly initialized reward model is trained
on the other four folds and evaluated on the held-out one. One implementation
detail is worth stating precisely, since it affects how ``training seed'' should
be interpreted. The \texttt{--seed} argument is passed to Python's
\texttt{random} module, which governs the fold-assignment shuffle. PyTorch's own
generator is not seeded anywhere in the cross-validation path
(\texttt{\_kfold\_cv}), so every source of stochasticity drawn from it is
uncontrolled: the randomly initialized reward head, the shuffling of training
batches by the \texttt{DataLoader}, and the dropout masks applied during
training. Changing \texttt{--seed} therefore changes the fold assignment, while
these other sources vary between process executions regardless of the seed value.
Together they explain why two executions launched with the same seed value
produced different results (Chapter~\ref{ch:results},
Section~\ref{sec:results-final-campaign}), and they are discussed as a source of
run-to-run variation in Chapter~\ref{ch:discussion},
Section~\ref{sec:disc-stochasticity}.

The final training configuration, used for every reported reward-model result, is:
\url{FacebookAI/roberta-base} backbone, 5 epochs, batch size 8, learning rate
$2\times10^{-5}$ with a linear warmup over 10\% of total steps, maximum sequence
length 512 tokens (article truncated to 2{,}000 characters before tokenization,
concatenated with the summary), and 5-fold cross-validation. This is confirmed by
\url{slurm/train_rm_1000.sh}, the Slurm script used for the headline
$n=1{,}000$ result. One inconsistency in the codebase is worth flagging explicitly,
since it could otherwise look like an error in the reported results: the in-code
default for \texttt{--backbone} in both \texttt{train.py} and
\texttt{run\_training.py} is \url{microsoft/deberta-v3-base}, a leftover from an
earlier version of the training script. Every Slurm script used to produce a
result reported in this thesis overrides this default and passes
\texttt{--backbone FacebookAI/roberta-base} explicitly; the \texttt{deberta-v3-base}
default was never actually used to produce a reported number. The one exception is
the very first Bradley--Terry result on the $n=495$ set
(Chapter~\ref{ch:results}, Section~\ref{sec:results-rm-baseline}), which predates
the standardization on \texttt{roberta-base} and is reported as such.

\paragraph{Source-document truncation.} The reward model does not see the full
source article. Two limits apply in sequence. The article string is first
truncated to its leading 2{,}000 characters
(\texttt{--max-article-chars}, applied in \texttt{PreferencePairDataset}), and the
truncated article is then concatenated with the candidate summary and tokenized
under a 512-token cap (\texttt{--max-length}), so the article prefix is shortened
further to make room for the summary. Subset selection admits articles up to
8{,}000 characters (Section~\ref{sec:impl-data}), so for the longer articles in
the pool the reward model observes well under half of the source text.

This matters for interpreting the reward-model results. The evaluator that
produced the preference labels scores summaries against the article under its own
chunking scheme, whereas the reward model must reconstruct those labels from a
prefix only. Where the evidence supporting or contradicting a claim lies beyond
the retained prefix, that evidence is simply unavailable to the reward model, and
the corresponding preference is not recoverable from its input in principle rather
than merely in practice. How large this effect is was not measured here, and the
near-chance accuracies reported in Chapter~\ref{ch:results} should not be
attributed to truncation without evidence; it is one plausible contributing factor
among several discussed in Chapter~\ref{ch:discussion},
Section~\ref{sec:disc-weak-signal}. Varying \texttt{--max-article-chars} while
holding everything else fixed would isolate it, and is recorded as future work.

\section{Compute Infrastructure}
\label{sec:impl-infra}

All experiments run on MOGON NHR (\texttt{mogon-nhr-01.zdv.uni-mainz.de}), the
University of Mainz's high-performance computing cluster, using Slurm partition
\texttt{a100dl} under account \texttt{nhr-haloed}, on NVIDIA A100-SXM4-40GB GPUs.
All Slurm scripts producing results reported in Chapter~\ref{ch:results} exclude
one cluster node found to have a faulty CUDA runtime, via
\texttt{\#SBATCH --exclude}; this affects scheduling only and not the computation
performed.

\section{Reproducibility}
\label{sec:impl-repro}

The pipeline is config-driven throughout: dataset, generation, and judging
parameters live in \url{configs/*.yaml} and can be overridden per run from the
command line, and every script accepts explicit flags for the parameters that
matter for a given experiment (margin, $\alpha$, AlignScore evaluation mode,
reward-model hyperparameters, and so on) rather than hardcoding them.

Seeding is partial, and the limits should be stated exactly. Subset selection and
cross-validation fold assignment are explicitly seeded through Python's
\texttt{random} module. PyTorch's generator is not seeded in the
cross-validation path, so weight initialization, batch ordering and dropout are
not fixed, and exact run-level reproduction from the command-line seed alone is
not guaranteed: two executions with identical arguments can differ. Every
pipeline script writes a JSON run-metadata file
(\url{src/utils/logging.py}, \texttt{save\_run\_metadata}) recording its
configuration, output artifact paths, and a timestamped run ID, so that any
reported number can in principle be traced back to the exact invocation that
produced it.

The configuration required to reproduce the $n=1{,}000$ result is given in full
in Table~\ref{tab:methodology-summary} (Chapter~\ref{ch:methodology}), together
with the six training seeds listed in Chapter~\ref{ch:results},
Table~\ref{tab:seed-validation}, and is implemented by
\url{slurm/train_rm_1000.sh} and \url{slurm/mode_nli_seed_confirm.sh}. The
codebase, including all scripts referenced in this chapter, is available at the
project's GitHub repository: \url{https://github.com/hasnat23/master-thesis-rlaif-gca}.

\chapter{Results}
\label{ch:results}

This chapter reports the experiments in the order they were carried out. Two
intermediate findings did not survive repetition: a hyperparameter configuration
that initially appeared promising (Section~\ref{sec:results-alpha0-validation})
and a two-run estimate that proved inconclusive
(Section~\ref{sec:results-mode-sweep}). Both are reported, since they motivated
the repeated-run protocol described in Chapter~\ref{ch:methodology},
Section~\ref{sec:methodology-rm-eval}, and omitting them would misrepresent how
much of the final estimate rests on distinguishing an effect from run-to-run
variance. Sections~\ref{sec:results-formula-opt} to~\ref{sec:results-mode-sweep}
are exploratory in the sense of Chapter~\ref{ch:methodology},
Section~\ref{sec:methodology-exploratory}; Section~\ref{sec:results-scale}
examines how the selected configuration behaves as the dataset grows.

\section{Early Prototype Results (Historical Context)}
\label{sec:results-early}

An earlier version of the pipeline produced results under a DPO-based design
before the study narrowed to reward-model comparison
(Chapter~\ref{ch:methodology}, Section~\ref{sec:methodology-evolution}). They are
reported here as context for that change, not as findings of the present study.

A FLAN-T5 baseline on a 20-sample pilot gave
ROUGE-1 0.3278, ROUGE-2 0.1201, ROUGE-L 0.2341, and BERTScore F1 0.8612. On a
200-sample subset (subset seed 42), the original AlignScore-based preference
construction, using the proposal-era margin of 0.05 and $\alpha = 0.5$, produced 154
usable holistic pairs out of 200 (77.0\%, mean scores 0.7324 for the low-temperature
candidate and 0.6530 for the high-temperature one) and 157 usable GCA pairs
(78.5\%, mean GCA scores 0.4066 and 0.3337). The two methods agreed on 121 of 200
decisions (60.5\%).

Mistral-7B-Instruct-v0.3 was then fine-tuned with DPO on both preference sets
(Chapter~\ref{ch:background}, Section~\ref{sec:bg-dpo}): training loss reached
0.6902 for the holistic condition and 0.6913 for the GCA condition, each in
roughly 134.5 seconds on an A100. On 200 held-out articles, baseline AlignScore was
0.7965, DPO-Holistic reached 0.8017, and DPO-GCA reached 0.7996. Only DPO-GCA's
ROUGE-L differed from baseline at the conventional threshold (Wilcoxon
$p=0.015$); the AlignScore differences between conditions were small and not
distinguishable from noise. These downstream results illustrate the attribution
problem described in Chapter~\ref{ch:methodology},
Section~\ref{sec:methodology-evolution}: differences of this magnitude cannot be
traced to the preference-construction procedure rather than to the optimization
stage, which is why the final study measures the preference data directly.

\section{Bradley--Terry Reward-Model Baseline}
\label{sec:results-rm-baseline}

The first Bradley--Terry reward-model result used the $n=495$ set (subset seed
100) with a \texttt{microsoft/deberta-v3-base} backbone, predating the
standardization on \texttt{FacebookAI/roberta-base}
(Chapter~\ref{ch:implementation}, Section~\ref{sec:impl-rm-training}). Under 5-fold
cross-validation, the holistic reward model reached 58.1\% mean pairwise accuracy
and the GCA reward model reached 54.6\%, both above the 50\% chance level.

The first result on the $n=1{,}000$ set, with the \texttt{roberta-base} backbone
and the original configuration ($\alpha = 0.5$, margin 0.05), is summarized in
Table~\ref{tab:rm-baseline-1000}. This is the baseline the rest of this chapter's
optimization campaign works from: holistic outperforms GCA by 1.2 percentage
points.

\begin{table}[htbp]
\centering
\caption[Bradley--Terry reward-model baseline at $n=1{,}000$]{Bradley--Terry reward-model baseline at $n=1{,}000$ ($\alpha=0.5$,
margin $0.05$, \texttt{roberta-base}).}
\label{tab:rm-baseline-1000}
\begin{tabular}{lrr}
\toprule
Condition & Mean pairwise accuracy & Gap (GCA $-$ Holistic) \\
\midrule
Holistic RM & 57.20\% & \multirow{2}{*}{$-1.20$ pp} \\
GCA RM      & 56.00\% & \\
\bottomrule
\end{tabular}
\end{table}

\section{GCA Formula Optimization}
\label{sec:results-formula-opt}

To understand why GCA underperformed, disagreement between the holistic and GCA
decisions at $\alpha = 0.5$ was analyzed directly on the $n=1{,}000$ pairs
(\texttt{analysis/gca\_vs\_holistic\_analysis.py}). The two methods disagreed on
17.8\% of pairs. GCA scores showed higher discrimination between candidates than
holistic scores (mean absolute score difference 0.2392 versus 0.1527), but only
82.2\% of GCA decisions agreed with the corresponding holistic decision, suggesting
that the extra discrimination was at least partly noise rather than signal.

Nine aggregation strategies were then compared by how well each agreed with the
holistic decision on the same 1{,}000 pairs
(\texttt{analysis/test\_gca\_formulas.py}, \texttt{analysis/test\_advanced\_aggregation.py}).
Table~\ref{tab:formula-comparison} reports the results. The simple mean
($\alpha = 0.0$) had the highest agreement of any strategy tested, 15.4 percentage
points above the original penalty formula.

\begin{table}[htbp]
\centering
\caption[Agreement with the holistic decision across nine aggregation strategies]{Agreement with the holistic decision across nine GCA aggregation
strategies, $n=1{,}000$ pairs.}
\label{tab:formula-comparison}
\begin{tabular}{lrr}
\toprule
Strategy & Agreement with holistic & vs.\ $\alpha=0.5$ baseline \\
\midrule
Simple mean ($\alpha=0.0$)   & 97.6\% & $+15.4$ pp \\
Confidence-weighted          & 96.0\% & $+14.0$ pp \\
Ensemble                     & 93.3\% & $+11.1$ pp \\
Length-weighted              & 91.5\% & $+9.3$ pp \\
Weighted by position         & 88.2\% & $+6.0$ pp \\
Robust mean (10\% trim)      & 85.4\% & $+3.2$ pp \\
Harmonic mean                & 82.9\% & $+0.7$ pp \\
Original penalty ($\alpha=0.5$) & 82.2\% & baseline \\
Minimum                      & 78.7\% & $-3.5$ pp \\
\bottomrule
\end{tabular}
\end{table}

\begin{figure}[htbp]
\centering
\begin{tikzpicture}
\begin{axis}[
  thesisplot,
  width=0.80\textwidth, height=6.2cm,
  xbar, bar width=8pt,
  xmin=70, xmax=100,
  xlabel={agreement with the holistic decision (\%)},
  ytick=data,
  symbolic y coords={Minimum,{Penalty $\alpha{=}0.5$},Harmonic mean,
                     Robust mean,Position-weighted,Length-weighted,
                     Ensemble,Confidence-weighted,{Simple mean $\alpha{=}0.0$}},
  y dir=reverse,
  nodes near coords, nodes near coords style={font=\scriptsize},
  every node near coord/.append style={/pgf/number format/fixed,
        /pgf/number format/fixed zerofill, /pgf/number format/precision=1},
]
\addplot[gcabar] coordinates {
  (78.7,Minimum) (82.2,{Penalty $\alpha{=}0.5$}) (82.9,Harmonic mean)
  (85.4,Robust mean) (88.2,Position-weighted) (91.5,Length-weighted)
  (93.3,Ensemble) (96.0,Confidence-weighted) (97.6,{Simple mean $\alpha{=}0.0$})};
\end{axis}
\end{tikzpicture}
\caption[Agreement across the nine aggregation strategies]{Agreement with the holistic decision across the nine aggregation
strategies of Table~\ref{tab:formula-comparison} ($n=1{,}000$ pairs). Agreement
measures consistency between two procedures driven by the same evaluator, not
external factual correctness.}
\label{fig:formula-agreement}
\end{figure}

On the strength of this result, the default $\alpha$ was changed from 0.5 to 0.0
throughout the codebase (\url{src/judging/gca.py},
\url{src/judging/build_reward_preferences.py}, and every Slurm script).

Two qualifications apply to this selection. First, the criterion used in
Table~\ref{tab:formula-comparison} is agreement with the holistic decision, which
is a measure of internal consistency between two procedures driven by the same
evaluator, not a measure of external factual correctness. A strategy could agree
closely with holistic scoring while reproducing the same errors, and nothing in
this comparison would distinguish those cases. The criterion is used here as a
diagnostic for whether the penalty term was amplifying disagreement, not as
evidence that the simple mean produces better labels. Second, this comparison and
the evaluator-mode sweep that follows are exploratory in the sense defined in
Chapter~\ref{ch:methodology}, Section~\ref{sec:methodology-exploratory}: the final
configuration was chosen after observing $n=1{,}000$ outcomes, so results on that
setting cannot also serve as independent confirmation of the choice. The
$n=5{,}000$ and $n=10{,}000$ experiments do not fully remedy this, because their
subsets are nested supersets of the $n=1{,}000$ set rather than independent
samples (Chapter~\ref{ch:methodology}, Section~\ref{sec:methodology-nesting}).

\section{Does the Formula Change Alone Flip the Result?}
\label{sec:results-alpha0-validation}

The formula change was then validated the same way the baseline was measured: by
training reward models on the resulting preferences and checking pairwise
accuracy, holding the judge mode fixed at the original \texttt{nli\_sp} setting.
Table~\ref{tab:alpha0-validation} shows that $\alpha = 0.0$ did not, on its own,
close the gap. Holistic accuracy improved more than GCA accuracy did, and the gap
between them widened slightly rather than closing.

\begin{table}[htbp]
\centering
\caption[Effect of the aggregation-formula change alone]{Effect of the formula change alone ($n=1{,}000$, judge mode
\texttt{nli\_sp} unchanged).}
\label{tab:alpha0-validation}
\begin{tabular}{lrrr}
\toprule
Metric & $\alpha=0.5$ (baseline) & $\alpha=0.0$ & Change \\
\midrule
Holistic mean accuracy & 0.572 & 0.586 & $+0.014$ \\
GCA mean accuracy      & 0.560 & 0.561 & $+0.001$ \\
Gap (Holistic $-$ GCA) & 0.012 & 0.025 & $+0.013$ (wider) \\
\bottomrule
\end{tabular}
\end{table}

A nine-configuration hyperparameter search was run next, sweeping learning rate
over $\{1\times10^{-5}, 2\times10^{-5}, 5\times10^{-5}\}$ and epoch count over
$\{3, 5, 7\}$, still at $\alpha = 0.0$. One configuration, learning rate
$1\times10^{-5}$ with 7 epochs, produced a GCA advantage of $+0.014$ (GCA 0.586
versus holistic 0.572). A same-seed confirmation rerun of that exact
configuration, however, produced the opposite result: holistic 0.577 versus GCA
0.560, a gap of $-0.017$. The apparent advantage did not reproduce. This result
is reported in full because of what it implies: at this dataset size, a single
training run, even with 5-fold cross-validation, can produce a result that looks
like a real effect but is actually within the range of ordinary fold-to-fold and
initialization noise. This directly motivated moving away from hyperparameter
tuning toward the judge-mode sweep below, and toward the multi-seed validation
protocol used for the headline result.

\section{AlignScore Judge-Mode Sweep}
\label{sec:results-mode-sweep}

With $\alpha = 0.0$ fixed, four AlignScore evaluation modes
(Chapter~\ref{ch:background}, Section~\ref{sec:bg-alignscore}) were compared under
the same $n=1{,}000$ training setup. Table~\ref{tab:mode-sweep} shows that the
choice of mode, not only the aggregation formula, is associated with whether GCA
leads at all. The \texttt{nli} mode produced the largest GCA
advantage of any configuration tested so far; \texttt{nli\_sp}, the mode used for
every result up to this point, is actually one of the two modes that favor
holistic.

\begin{table}[htbp]
\centering
\caption[AlignScore evaluator-mode sweep]{AlignScore judge-mode sweep ($n=1{,}000$, $\alpha=0.0$).}
\label{tab:mode-sweep}
\begin{tabular}{lrrr}
\toprule
Mode & Holistic mean acc.\ & GCA mean acc.\ & Gap (GCA $-$ Holistic) \\
\midrule
\texttt{nli}    & 0.523 & 0.583 & $+0.060$ \\
\texttt{bin}    & 0.510 & 0.564 & $+0.054$ \\
\texttt{bin\_sp} & 0.554 & 0.562 & $+0.008$ \\
\texttt{nli\_sp} & 0.578 & 0.557 & $-0.021$ \\
\bottomrule
\end{tabular}
\end{table}

Figure~\ref{fig:mode-sweep} plots the same values grouped by whether the
evaluator performs its own internal splitting.

\begin{figure}[htbp]
\centering
\begin{tikzpicture}
\begin{axis}[
  thesisplot,
  width=0.82\textwidth, height=5.6cm,
  ybar, bar width=13pt,
  ymin=0.48, ymax=0.62,
  ylabel={mean pairwise accuracy},
  symbolic x coords={nli,bin,bin\_sp,nli\_sp},
  xtick=data,
  legend style={at={(0.5,-0.22)}, anchor=north, legend columns=2},
  nodes near coords, nodes near coords style={font=\scriptsize,
        /pgf/number format/fixed, /pgf/number format/precision=3},
]
\addplot[holbar] coordinates {(nli,0.523) (bin,0.510) (bin\_sp,0.554) (nli\_sp,0.578)};
\addplot[gcabar] coordinates {(nli,0.583) (bin,0.564) (bin\_sp,0.562) (nli\_sp,0.557)};
\draw[dashed, gray!70] ({rel axis cs:0,0}|-{axis cs:nli,0.5}) --
                       ({rel axis cs:1,0}|-{axis cs:nli,0.5});
\node[font=\scriptsize, gray!60!black, anchor=south east]
      at ({rel axis cs:0.995,0}|-{axis cs:nli,0.5}) {chance};
\legend{Holistic, GCA}
\end{axis}
\end{tikzpicture}
\caption[Evaluator-mode sweep]{Evaluator-mode sweep ($n=1{,}000$, $\alpha=0.0$). The two left-hand
modes do not split the input internally; the two right-hand modes do. Each bar is
a single run, so differences of this size fall within the run-to-run variation
documented in Section~\ref{sec:results-final-campaign}, and the comparison is
exploratory.}
\label{fig:mode-sweep}
\end{figure}

The ordering in Table~\ref{tab:mode-sweep} follows a systematic pattern rather
than varying arbitrarily. The two modes in which GCA leads by a clear margin,
\texttt{nli} and
\texttt{bin}, are precisely the two that operate on the context and claim without
splitting them. The two modes in which the advantage is negligible or negative,
\texttt{bin\_sp} and \texttt{nli\_sp}, are the split variants, in which the
evaluator itself segments the input and aggregates over the segments
(Chapter~\ref{ch:background}, Section~\ref{sec:bg-alignscore}). Holistic accuracy
also rises from 0.510--0.523 in the unsplit modes to 0.554--0.578 in the split
modes, which is consistent with the split modes supplying a stronger holistic
signal rather than with GCA becoming worse.

The pattern suggests that external sentence-level decomposition may contribute
most where the evaluator does not already decompose internally, and little where
it does; on this reading GCA and the evaluator's own splitting act partly as
substitutes. This sweep is exploratory and each mode was run once, so the
comparison across modes should be read as suggestive rather than established:
Section~\ref{sec:results-final-campaign} shows that single-run differences of this
magnitude fall within the range of run-to-run variation, and only the
\texttt{nli} configuration was subsequently validated across repeated runs.
Chapter~\ref{ch:discussion}, Section~\ref{sec:disc-interpretation} develops the
interpretation and states what would be required to test it.

Because a single run had already proved insufficient in
Section~\ref{sec:results-alpha0-validation}, the \texttt{nli} result was rerun at
the same seed value before any conclusion was drawn. The rerun gave holistic 0.543
and GCA 0.556, a gap of $+0.013$: smaller than the original $+0.060$, but in the
same direction. Over the fold-level differences from both runs, the 95\% interval
was $[+0.003, +0.064]$ and the corresponding Wilcoxon $p$-value was 0.084. Under
the independence caveat discussed above, neither figure was treated as decisive,
and further runs were carried out rather than concluding from two.

\section{Repeated Runs of the Selected Configuration (\texorpdfstring{$n=1{,}000$}{n=1,000})}
\label{sec:results-final-campaign}

The \texttt{nli} mode, $\alpha=0.0$, margin 0 configuration was fixed and run
across four further seed values, in addition to the original sweep and its
confirmation rerun, for six runs and 30 cross-validation folds in total.
Table~\ref{tab:seed-validation} lists every run. Note that the six runs span five
distinct seed values, not six: seed 42 appears twice, as the original sweep and
its rerun. As discussed in Chapter~\ref{ch:methodology},
Section~\ref{sec:methodology-rm-eval}, those two runs are not duplicates, because
the seed argument controls only fold assignment and not the reward model's weight
initialization.

\begin{table}[htbp]
\centering
\caption[Per-run results for the selected configuration]{Per-run results for the selected \texttt{nli}-mode configuration
($n=1{,}000$).}
\label{tab:seed-validation}
\begin{tabular}{lrrrr}
\toprule
Run & Seed & Holistic & GCA & Gap (GCA $-$ Holistic) \\
\midrule
Original sweep       & 42   & 0.523 & 0.583 & $+0.060$ \\
Confirmation          & 42   & 0.543 & 0.556 & $+0.013$ \\
Seed validation 1    & 7    & 0.556 & 0.546 & $-0.010$ \\
Seed validation 2    & 100  & 0.510 & 0.561 & $+0.051$ \\
Seed validation 3    & 314  & 0.520 & 0.552 & $+0.032$ \\
Seed validation 4    & 2026 & 0.525 & 0.564 & $+0.039$ \\
\bottomrule
\end{tabular}
\end{table}

GCA leads in five of the six runs; seed 7 is the only run in which holistic is
ahead. Figure~\ref{fig:per-run-gaps} plots the per-run gaps against the pooled
mean, which makes the spread easier to judge than the table alone: the individual
runs range from $-0.010$ to $+0.060$, a spread wide enough that any single run,
taken on its own, would have been weak evidence in either direction. Averaging
over all 30 fold-level readings gives the central estimate of this thesis, shown
in Table~\ref{tab:pooled-headline}: a mean gap of $+3.08$ percentage points.

\begin{figure}[htbp]
\centering
\begin{tikzpicture}[font=\footnotesize]
  % axes
  \def\xz{0}      % x origin
  \def\xs{95}     % x scale (per 0.01 gap) -- gaps range -0.01..0.06
  \draw[->, thick] (-1.3,0) -- (6.9,0) node[right] {gap};
  \foreach \g/\lab in {-0.01/-0.01, 0/0, 0.02/+0.02, 0.04/+0.04, 0.06/+0.06} {
    \draw ({\g*\xs},0.12) -- ({\g*\xs},-0.12) node[below] {\lab};
  }
  % zero reference line
  \draw[dashed, gray] (0,-0.4) -- (0,5.4);
  % pooled mean line
  \draw[thick, blue!60!black, dashed] ({0.0308*\xs},-0.4) -- ({0.0308*\xs},5.4);
  \node[blue!60!black, anchor=south, align=center] at ({0.0308*\xs},5.4)
    {pooled mean\\$+0.0308$};
  % runs (y from top to bottom)
  \foreach \y/\seed/\gap in {
      5/42/0.060,
      4.2/42/0.013,
      3.4/7/-0.010,
      2.6/100/0.051,
      1.8/314/0.032,
      1.0/2026/0.039} {
    \node[anchor=east] at (-1.35,\y) {seed \seed};
    \fill[orange!80!black] ({\gap*\xs},\y) circle (2.4pt);
    \draw[gray!60] (0,\y) -- ({\gap*\xs},\y);
  }
\end{tikzpicture}
\caption[Per-run GCA$-$Holistic accuracy gap at $n=1{,}000$]{Per-run GCA$-$Holistic accuracy gap at $n=1{,}000$, against the pooled
mean. Five of six runs favor GCA, but the spread across runs is wide relative to
the pooled effect, which is why no single run was treated as conclusive.}
\label{fig:per-run-gaps}
\end{figure}

\begin{table}[htbp]
\centering
\caption[Averaged result across six runs at $n=1{,}000$]{Averaged result across 6 runs / 30 cross-validation folds at
$n=1{,}000$. The interval and $p$-value are computed over folds, which are not
independent observations, and are therefore exploratory and anti-conservative
(Chapter~\ref{ch:methodology}, Section~\ref{sec:methodology-rm-eval}).}
\label{tab:pooled-headline}
\begin{tabular}{rrrrr}
\toprule
Holistic mean & GCA mean & Gap & 95\% fold-level CI & Fold-level Wilcoxon $p$ \\
\midrule
0.5295 & 0.5603 & \textbf{+0.0308} & $[+0.013, +0.047]$ & 0.0034 \\
\bottomrule
\end{tabular}
\end{table}

At the fold level, GCA leads in 22 of the 30 folds (73\%), with no ties. The
interval and $p$-value in Table~\ref{tab:pooled-headline} are computed as
described in Chapter~\ref{ch:background}, Section~\ref{sec:bg-stats}, using
10{,}000 bootstrap resamples with a fixed NumPy seed of 42.

These uncertainty estimates require an explicit caveat and are not presented as
confirmatory. The 30 folds are not 30 independent observations: the five folds
within a run partition one dataset, so their training sets overlap heavily and
their errors are correlated. A procedure that assumes independence therefore
treats the evidence as stronger than it is, and the quoted $p$-value is
anti-conservative. What the fold-level analysis supports is that the difference is
consistent in direction across runs and stable in magnitude, not that it meets a
particular significance threshold. Chapter~\ref{ch:discussion},
Section~\ref{sec:disc-stats-caveat} examines what a more conservative treatment
would and would not permit.

A second property should be read alongside the gap. Both conditions sit close to
chance in absolute terms: 0.5295 and 0.5603 are 2.95 and 6.03 percentage points
above the 0.5 chance level. The comparison is therefore between two weak
preference signals, and the finding concerns which is more consistently
recoverable, not whether either yields a strong reward model
(Section~\ref{sec:disc-weak-signal}).

Table~\ref{tab:evidence-progression} tracks how the estimate changed as runs were
added. The interval narrows at each stage while the point estimate stays near
three percentage points. Stability of the effect size as precision improves is
consistent with a real underlying difference rather than with an artifact that
grows as data accumulates, though the same independence caveat applies to every
row of the table.

\begin{table}[htbp]
\centering
\caption[Progression of the pooled estimate as runs were added]{Progression of the pooled statistical estimate as runs were added. The
$p$-values are subject to the independence caveat discussed in
Chapter~\ref{ch:discussion}, Section~\ref{sec:disc-stats-caveat}.}
\label{tab:evidence-progression}
\begin{tabular}{lrrrrr}
\toprule
Stage & Runs & Folds & Pooled gap & 95\% CI & Wilcoxon $p$ \\
\midrule
Two-run  & 2 & 10 & $+0.034$ & $[+0.003, +0.064]$ & 0.084 \\
Five-run & 5 & 25 & $+0.029$ & $[+0.009, +0.048]$ & 0.0123 \\
Six-run  & 6 & 30 & $+0.031$ & $[+0.013, +0.047]$ & 0.0034 \\
\bottomrule
\end{tabular}
\end{table}

\section{Scale Robustness: 5{,}000 and 10{,}000 Samples}
\label{sec:results-scale}

The selected configuration was then run once at $n=5{,}000$ and once at
$n=10{,}000$. These subsets are nested supersets of the $n=1{,}000$ set rather
than independent samples (Chapter~\ref{ch:methodology},
Section~\ref{sec:methodology-nesting}), so they measure how the effect behaves as
data is added to the same pool rather than providing out-of-sample replication.
Unlike the
$n=1{,}000$ result, each of these runs is a single run, not a six-run pool, for
computational-budget reasons stated in Chapter~\ref{ch:methodology},
Section~\ref{sec:methodology-rm-eval}; the comparison in
Table~\ref{tab:scale-robustness} should be read with that difference in mind
rather than treated as three results of equal statistical weight.

\begin{table}[htbp]
\centering
\caption[Selected configuration at three dataset sizes]{Selected configuration at three dataset sizes. Only the $n=1{,}000$ row
is averaged over repeated runs; the other two are single runs on nested datasets.}
\label{tab:scale-robustness}
\begin{tabular}{rrrr}
\toprule
Dataset size & Holistic mean & GCA mean & Gap (GCA $-$ Holistic) \\
\midrule
1{,}000  & 0.5295 & 0.5603 & $+0.0308$ \\
5{,}000  & 0.5788 & 0.5746 & $-0.0042$ \\
10{,}000 & 0.5827 & 0.5862 & $+0.0035$ \\
\bottomrule
\end{tabular}
\end{table}

\begin{figure}[htbp]
\centering
\begin{tikzpicture}
\begin{axis}[
  thesisplot,
  width=0.82\textwidth, height=5.6cm,
  ybar, bar width=15pt,
  ymin=0.48, ymax=0.63,
  ylabel={mean pairwise accuracy},
  xlabel={preference pairs},
  symbolic x coords={1{,}000,5{,}000,10{,}000},
  xtick=data,
  legend style={at={(0.03,0.97)}, anchor=north west, legend columns=1},
]
\addplot[holbar] coordinates {(1{,}000,0.5295) (5{,}000,0.5788) (10{,}000,0.5827)};
\addplot[gcabar] coordinates {(1{,}000,0.5603) (5{,}000,0.5746) (10{,}000,0.5862)};
\draw[dashed, gray!70] ({rel axis cs:0,0}|-{axis cs:1{,}000,0.5}) --
                       ({rel axis cs:1,0}|-{axis cs:1{,}000,0.5});
\node[font=\scriptsize, gray!60!black, anchor=south west]
      at ({rel axis cs:0.30,0}|-{axis cs:1{,}000,0.5}) {chance};
\legend{Holistic, GCA}
\end{axis}
\end{tikzpicture}
\caption[Accuracy at the three dataset sizes]{Accuracy at the three dataset sizes. Absolute accuracy rises for both
conditions as data is added, while the difference between them contracts. Only
the $n=1{,}000$ bars are averaged over repeated runs; the other two are single
runs on datasets that are nested supersets of the $n=1{,}000$ set
(Chapter~\ref{ch:methodology}, Section~\ref{sec:methodology-nesting}).}
\label{fig:scale-results}
\end{figure}

The $n=1{,}000$ advantage is not reproduced with comparable magnitude at
$n=5{,}000$: holistic performs marginally better there, though the gap is small.
At $n=10{,}000$ the direction returns in GCA's favor, but the magnitude is an
order of magnitude smaller than at $n=1{,}000$. Neither larger-scale result was
repeated across runs the way the $n=1{,}000$ result was, so it cannot be ruled
out that some of this movement reflects the same run-to-run variation documented
in Section~\ref{sec:results-alpha0-validation} rather than a change in the
underlying effect as dataset size grows. What can be said without qualification is
that a single dataset size is not sufficient to characterize the difference
between the two conditions, and that reporting only $n=1{,}000$ would have
overstated how consistently the GCA advantage is observed.

\section{Summary of Results}
\label{sec:results-summary}

Table~\ref{tab:rq-summary} summarizes how the results in this chapter bear on the
research questions from Chapter~\ref{ch:methodology}, Section~\ref{sec:methodology-rqs}.
Full interpretation is left to Chapter~\ref{ch:discussion}.

\begin{table}[htbp]
\centering
\caption[Summary of findings against the research questions]{Summary of findings against the research questions.}
\label{tab:rq-summary}
\renewcommand{\arraystretch}{1.25}
\begin{tabularx}{\textwidth}{@{}>{\bfseries}l X@{}}
\toprule
& Finding \\
\midrule
RQ1 &
  At $n=1{,}000$, GCA leads in 5 of 6 runs with a mean difference of
  $+3.08$ percentage points. The run-level two-sided test does not meet
  $p<0.05$, and the larger single-run experiments do not reproduce an advantage
  of comparable magnitude. \\
RQ2 &
  Aggregation strategy and evaluator mode are associated with substantial
  outcome differences. The evaluator-mode comparison is exploratory, because
  every mode except the subsequently selected \texttt{nli} configuration was
  evaluated only once. \\
RQ3 &
  The direction is relatively stable across the six $n=1{,}000$ runs, but is not
  reproduced at comparable magnitude in the single $n=5{,}000$ and $n=10{,}000$
  runs, which use nested datasets. \\
\bottomrule
\end{tabularx}
\end{table}

\chapter{Discussion}
\label{ch:discussion}

\section{What the Experiments Establish, and What They Do Not}
\label{sec:disc-claims}

The experiments support a narrow and conditional claim. Under one specific
configuration, at one dataset size, preferences constructed by sentence-level
aggregation are more consistently recoverable by a Bradley--Terry reward model
than preferences constructed holistically from the same evaluator and the same
candidate summaries. The difference averages roughly three percentage points over
six repeated runs and is consistent in direction across five of them.

Four claims that might appear to follow do not. The results do not establish that
GCA labels are factually more accurate, since no independent annotation was
collected and AlignScore supplies both sets of labels
(Section~\ref{sec:disc-reliability}). They do not establish that a summarization
model trained on GCA preferences produces more faithful summaries, since the final
study contains no policy-optimization stage. They do not establish that GCA is
generally preferable, since Chapter~\ref{ch:results} reports configurations and
scales at which it is not. And they do not establish a strong effect in absolute
terms, since both reward models remain near chance
(Section~\ref{sec:disc-weak-signal}). What the study contributes is evidence about
\emph{when} feedback granularity changes the structure of a preference-learning
signal, and the remainder of this chapter examines the conditions under which it
does.

\section{Why Granularity Interacts With Evaluator Design}
\label{sec:disc-interpretation}

The advantage reported in Chapter~\ref{ch:results},
Section~\ref{sec:results-final-campaign} is not a property of GCA in the abstract.
It appears only under a specific combination of two design choices,
simple-mean aggregation ($\alpha=0.0$) and the AlignScore \texttt{nli} evaluation
mode, and Section~\ref{sec:results-alpha0-validation} already showed that the
aggregation-formula change alone is not sufficient: at $\alpha=0.0$ under the
\texttt{nli\_sp} mode, holistic still led. This section asks why the
judge mode specifically makes the difference.

AlignScore's implementation \citep{zha2023alignscore} suggests a plausible
mechanism, though the present experiments do not test it directly. The \texttt{\_sp} suffix in \texttt{nli\_sp} enables
a preprocessing and aggregation step that is internal to AlignScore itself: the
context is split into roughly 350-token chunks and the claim is split into
sentences, and the final score is computed as, in the library's own
\texttt{inference\_per\_example} function, the maximum entailment score across
context chunks for each claim sentence, averaged across claim sentences
(\texttt{.max(dim=0).values.mean()}). Plain \texttt{nli} mode skips this
entirely and scores the full context against the full claim in one pass. The
chunk-then-aggregate pattern is not unique to AlignScore: SummaC applies the same
idea, segmenting document and summary and reducing a matrix of sentence-level
entailment scores to a single figure \citep{laban2022summac}. What differs
between evaluators is where that decomposition happens, which is precisely the
distinction this section turns on.

This has a direct consequence for what the holistic and GCA conditions actually
compute under each mode. Under \texttt{nli\_sp}, a holistic call already
decomposes the candidate summary into sentences internally, scores each one
against the best-matching chunk of the article, and averages the results, which
is functionally close to what GCA does \emph{externally} at $\alpha=0.0$: score
each summary sentence against the article and take the mean. When GCA is applied
on top of a judge that is already doing something close to sentence-level
decomposition and mean-pooling internally, its external decomposition adds little
genuinely new information, which is consistent with holistic and GCA performing
similarly, and even favoring holistic, under \texttt{nli\_sp}
(Table~\ref{tab:mode-sweep}). Under plain \texttt{nli}, by contrast, a holistic
call scores the entire summary against the (length-truncated) article as one
atomic unit, with no internal sentence decomposition at all. Here, GCA's external
per-sentence scoring is the only source of sentence-level localization in the
pipeline, and it is exactly under this mode that GCA shows a measurable
advantage.

The stored $n=200$ preference data allow one component of this account to be
checked directly rather than inferred. If a holistic \texttt{nli\_sp} call
internally splits the claim into sentences and averages, then the holistic score
of a summary should closely track the mean of the separately computed
per-sentence scores of that same summary. Across all 400 summaries in that set
the two quantities correlate at $r = 0.983$, agree to within $0.001$ for 234 of
400 summaries, and have a median absolute difference of $0.000$
(Figure~\ref{fig:mean-vs-holistic}). Under \texttt{nli\_sp}, holistic scoring and
mean-aggregated sentence scoring are therefore close to the same computation,
which is consistent with the observation that GCA at $\alpha=0.0$ adds little
under that mode.

\begin{figure}[htbp]
\centering
\begin{tikzpicture}
\begin{axis}[
  thesisplot,
  width=0.62\textwidth, height=6.2cm,
  xlabel={mean of per-sentence scores},
  ylabel={holistic score},
  xmin=0, xmax=1, ymin=0, ymax=1,
  axis equal image, grid=major,
]
\addplot[only marks, mark=*, mark size=1.0pt,
         draw=orange!70!black, fill=orange!55, fill opacity=0.55]
        table[x=meansent, y=holistic] {figures/mean_vs_holistic.dat};
\addplot[domain=0:1, samples=2, dashed, gray!70, line width=0.7pt] {x};
\end{axis}
\end{tikzpicture}
\caption[Holistic score vs.\ mean of per-sentence scores]{Holistic AlignScore against the mean of the separately computed
per-sentence scores, for all 400 candidate summaries in the $n=200$ preference
set under evaluator mode \texttt{nli\_sp}. The dashed line is $y=x$. The two
quantities are near-identical ($r = 0.983$), which is what the split mode's
internal sentence-averaging predicts. The corresponding files for the final
\texttt{nli} configuration are not retained in the repository, so the same check
could not be run for the unsplit mode.}
\label{fig:mean-vs-holistic}
\end{figure}

This measurement supports the part of the account concerning \texttt{nli\_sp},
but the account as a whole remains an interpretation of an exploratory sweep, and
three limitations constrain it. The mode comparison rests on a
single run per mode, and Section~\ref{sec:disc-stochasticity} shows that
single-run differences of this size fall within observed run-to-run variation;
only the selected \texttt{nli} configuration received repeated-run validation. No
ablation isolates the components of the \texttt{\_sp} preprocessing, such as chunk
size, from the mode change as a whole. And nothing here verifies empirically that
\texttt{nli\_sp}'s internal aggregation and GCA's external aggregation produce
comparable scores rather than merely comparable rankings. The proposition that
evaluator-internal decomposition reduces the marginal benefit of external GCA is
therefore offered as the most plausible reading of the observed pattern, and as a
hypothesis worth testing directly, rather than as a finding of this thesis
(Chapter~\ref{ch:conclusion}).

\section{How Much Weight the Statistical Evidence Can Carry}
\label{sec:disc-stats-caveat}

The $n=1{,}000$ result was first assessed from a fold-level interval and a
Wilcoxon $p$-value of 0.0034, computed over 30 pooled fold-level differences from
six runs. It is necessary to state plainly how much that figure could support, and
then to describe what changed once the campaign was extended.

Those 30 differences are not 30 independent observations. Within a single run the
five cross-validation folds partition one dataset: any two folds share four fifths
of their training data, and all five evaluate models trained from the same pool of
preference pairs. Their errors are correlated, so a test assuming independence
treats the evidence as stronger than it is. The effective sample size lies between
the number of runs and the number of folds, and closer to the former. The
fold-level $p$-value is therefore anti-conservative and is reported as an
exploratory summary, not as a hypothesis test, at either campaign size.

The appropriate unit of independent replication is the run. Applying the two-sided
Wilcoxon signed-rank test to the six run-level differences from the original
campaign ($+0.060$, $+0.013$, $-0.010$, $+0.051$, $+0.032$, $+0.039$) gives
$p = 0.0625$, which does not meet the conventional $0.05$ threshold. With six
observations, five of them positive, $0.0625$ is in fact the smallest two-sided
value the test can return, so that outcome reflected the number of runs available
at least as much as the size of the effect: no result with five wins and one loss
out of six paired observations can reach significance, regardless of how large or
small the true difference is.

This is a genuine limitation of a six-run design, not a property of the
underlying effect, and Chapter~\ref{ch:results},
Section~\ref{sec:results-extended-campaign} reports what happened when it was
addressed directly: the campaign was extended to twenty independently seeded
runs under a corrected, fully deterministic training procedure. GCA led in all
twenty, with a mean gap of $+3.95$ percentage points, a bootstrap 95\% confidence
interval of $[+3.19, +4.70]$ points that does not cross zero, and a run-level
two-sided Wilcoxon $p < 0.001$. This does meet the conventional significance
threshold, and the effect size is close to the original six-run estimate,
indicating that the earlier campaign's direction was not an artifact of having
too few runs to trust, only too few to prove.

Reaching significance changes what can be claimed, but it does not remove the
need for the qualifications that surrounded the six-run result. Statistical
significance at the run level establishes that GCA preferences are
\emph{consistently more learnable} than holistic preferences by this
reward-model architecture, under this evaluator configuration, at this dataset
size, to a degree unlikely to arise from run-to-run noise alone. It does not by
itself establish a large effect (Section~\ref{sec:disc-weak-signal} addresses
the absolute accuracy level), an effect that generalizes to other scales
(Section~\ref{sec:disc-scale}), or a mechanism (Section~\ref{sec:disc-interpretation}
remains an interpretation, not confirmed by the extended campaign). What the
evidence now supports is a directionally consistent, statistically significant
advantage of approximately three to four percentage points under this
configuration --- a materially stronger claim than the six-run campaign could
support, but one that is still scoped to the conditions under which it was
measured.

\section{Both Reward Models Are Weak in Absolute Terms}
\label{sec:disc-weak-signal}

A comparison between two conditions can be sound while both conditions perform
poorly, and that is the situation here. At $n=1{,}000$, the holistic reward model
reaches 0.5295 and the GCA reward model 0.5603 pairwise accuracy, against a chance
level of 0.5. In absolute terms these models recover very little of the preference
signal: the better of the two is correct on roughly 56 of every 100 held-out pairs,
where guessing would be correct on 50.

This materially qualifies what the thesis shows. The finding is that GCA
preferences are \emph{more learnable} than holistic preferences under the selected
configuration, not that either produces a reward model one would want to deploy.
A reward model at 0.56 accuracy would provide a very noisy training signal to any
downstream policy, and the significant $+3.95$ percentage-point advantage
established in the extended campaign (Chapter~\ref{ch:results},
Section~\ref{sec:results-extended-campaign}) is an improvement to a weak signal
rather than the difference between a weak and a strong one.

Several factors plausibly contribute to the low absolute numbers, and this thesis
can separate them only partially. The task itself is genuinely hard: both
candidate summaries in a pair are generated by the same model from the same
article and differ only in sampling temperature, so many pairs are near-ties in
factual quality, and a judge's preference between them may encode little
learnable structure. The RoBERTa-base reward model \citep{liu2019roberta} is also
small, and articles are truncated to 2{,}000 characters before tokenization
(Chapter~\ref{ch:implementation}, Section~\ref{sec:impl-rm-training}), so the model
may simply not see enough of the source to verify many claims. The accuracies at
$n=5{,}000$ and $n=10{,}000$ are higher for both conditions (0.58--0.59), which is
consistent with more training data helping the reward model in general, even as
the gap between conditions closes.

A preliminary ablation tested the truncation explanation directly by varying the
retained article length under the RoBERTa-base backbone, three seeds per setting
(Table~\ref{tab:truncation-ablation}). The result runs against the truncation
hypothesis as originally stated: accuracy was \emph{higher} with less retained
context (500 and 1{,}000 characters) than at the thesis's own 2{,}000-character
setting, for both conditions. A plausible reading is that 2{,}000 characters is
not actually more informative in practice, because the concatenation of article
prefix and summary already exceeds the 512-token cap at that length
(Chapter~\ref{ch:implementation}, Section~\ref{sec:impl-rm-training}), so how much
of the 2{,}000-character budget survives tokenization varies unpredictably with
summary length across examples. A shorter, more consistently-applied budget may
therefore produce a more uniform input distribution even though it discards more
raw text. This ablation does not yet isolate that explanation from a simple
backbone effect, since it was run only under RoBERTa-base; the two prepared but
unexecuted arms of the same ablation, under \texttt{deberta-v3-base} at 512 and
1{,}024 tokens, would hold the backbone fixed while varying context length and
are recorded as future work (Chapter~\ref{ch:conclusion}). Taken as it stands,
this result reframes rather than resolves the truncation question: raw article
length is not the limiting factor it was hypothesized to be, at least not in the
straightforward way of "more characters, more evidence, higher accuracy," and the
low absolute ceiling remains only partially explained.

\begin{table}[htbp]
\centering
\caption[Truncation ablation, RoBERTa-base arms]{Truncation ablation, RoBERTa-base arms (three seeds each). All three
rows use the 512-token cap; \texttt{max\_article\_chars} is the character budget
before tokenization. The \texttt{deberta-v3-base} arms (512 and 1{,}024 tokens)
were prepared but not completed; see Chapter~\ref{ch:conclusion}.}
\label{tab:truncation-ablation}
\begin{tabular}{rrrrr}
\toprule
Max article chars & Holistic & GCA & Gap (GCA $-$ Holistic) \\
\midrule
500   & 0.5553 & 0.5850 & +2.97 \\
1{,}000 & 0.5663 & 0.5903 & +2.40 \\
2{,}000 & 0.5290 & 0.5480 & +1.90 \\
\bottomrule
\end{tabular}
\end{table}

\section{Interpreting the Larger-Scale Results}
\label{sec:disc-scale}

The $n=1{,}000$ advantage was not reproduced with comparable magnitude in the
single runs at $n=5{,}000$ and $n=10{,}000$ (Chapter~\ref{ch:results},
Section~\ref{sec:results-scale}). Three things must be kept apart here: the
observed numerical behaviour, run-to-run variation, and hypotheses about
dataset-size effects. The observation is that two single runs produced gaps of
$-0.42$ and $+0.35$ percentage points where the pooled $n=1{,}000$ estimate was
$+3.08$. Whether this reflects a dataset-size effect is a separate question, and
the design does not answer it.

The best-supported explanation is statistical rather than substantive. The
$n=1{,}000$ estimate itself only stabilized after six pooled runs; the individual
per-run gaps in Table~\ref{tab:seed-validation}
ranged from $-0.010$ to $+0.060$, a spread wide enough that a single seed could
easily have landed on either side of zero. The $n=5{,}000$ and $n=10{,}000$
settings each contain a single run. Given how much the $n=1{,}000$ estimate moved
between its first two executions alone (from $+0.060$ down to $+0.013$), a single
run at a different dataset size landing outside the $n=1{,}000$ interval is not,
by itself, strong evidence that the underlying effect changes with dataset size.
It is at least as consistent with the same run-to-run variation already documented
at $n=1{,}000$, which simply went unmeasured at the larger sizes because only one
run was performed at each.

A substantive explanation is also possible, but this thesis does not have
evidence to support or rule it out: larger candidate pools could shift the
distribution of article lengths, sentence counts, or segmentation quality in ways
that change how much signal GCA's per-sentence scoring adds relative to holistic
scoring. This is stated here as an open hypothesis, not a finding, since nothing
in Chapter~\ref{ch:results} isolates article length, sentence count, or
segmentation quality as a variable.

Because these two explanations make different predictions, they are also
distinguishable in principle: if the scale effect is genuinely statistical noise,
repeated runs at $n=5{,}000$ should produce an interval overlapping the
$n=1{,}000$ result; if it reflects a real, scale-dependent change, they should
not. Repeating the $n=5{,}000$ experiment across several runs is the most direct
way to resolve this and is listed as future work in
Chapter~\ref{ch:conclusion}.

One hypothesis would predict this pattern, and it is worth stating even though the
present data cannot test it. At $n=1{,}000$ the reward
model is fitting a small dataset, and how cleanly the preference relation is
structured matters a great deal: a labeling that is internally more consistent is
easier to recover from few examples. As the dataset grows, the model sees more
examples of the same underlying relation, and additional data can compensate for a
noisier labeling. Under that account, the benefit of a cleaner preference
construction should be largest where data is scarce and should shrink as data
accumulates, which is the pattern observed. Consistent with this, absolute
accuracy rises for \emph{both} conditions between $n=1{,}000$ and $n=10{,}000$
(from 0.5295 and 0.5603 to 0.5827 and 0.5862), while the difference between them
contracts. The evidence is compatible with this explanation but does not establish
it, since scale and dataset composition vary together across these runs.

Read this way, the larger-scale results are informative rather than merely
disappointing. The observed pattern is consistent with the hypothesis that GCA may
provide a larger learnability advantage when preference data are limited, which
would also be the regime in which automatic preference construction is most
attractive, since it is where collecting human labels would be most burdensome.
The present design cannot establish such a boundary, because the larger-scale
settings contain single runs on nested datasets. What the study does contribute
here is the observation itself, together with a specification of the experiment
that would test it.

\section{Implications for RLAIF Preference Construction}
\label{sec:disc-implications}

Three implications follow for anyone constructing AI-generated preference data
from a fixed factuality judge, independent of whether the specific GCA method
studied here is adopted.

First, the aggregation formula used to combine sentence-level scores into a
preference label is not a minor implementation detail. The original
weakest-link-penalty formula ($\alpha=0.5$) actively hurt GCA relative to a plain
mean (Chapter~\ref{ch:results}, Table~\ref{tab:formula-comparison}), which runs
counter to the intuition that motivated it: a formula designed to be more
sensitive to localized errors was, in practice, more sensitive to noise than to
signal. Aggregation-formula choices of this kind should be validated
empirically rather than assumed correct from first principles. This is a
qualification worth attaching to the broader case for fine-grained supervision
\citep{wu2023finegrained}: locating errors more precisely helps only if the rule
that converts those locations into a training signal is itself sound.

Second, judge configuration can matter more than the granularity decision that
is nominally the point of comparison. Section~\ref{sec:disc-interpretation}'s
reading of AlignScore's mode behavior suggests that part of what looked like a
question about holistic-versus-sentence-level supervision was really a question
about how much sentence-level decomposition the judge was already doing
internally. Anyone comparing supervision granularities on top of a judge with
its own internal aggregation behavior should check what that judge already does
before attributing an effect to the granularity of their own pipeline. Surveys of
automatic judging emphasise sensitivity to prompt and configuration choices
\citep{zheng2023judging, gu2024survey}; the present result adds that a judge's
internal decomposition can also absorb the very manipulation an experiment
intends to study.

Third, a result validated at one dataset scale should not be treated as general
without a scale check. This is a caution this thesis can make credibly only
because Section~\ref{sec:disc-scale}'s robustness check was actually run rather
than assumed unnecessary; had the study stopped at the $n=1{,}000$ result, it
would have reported a materially less qualified conclusion than the evidence
supports.

\section{Judge Reliability Without Human Auditing}
\label{sec:disc-reliability}

No independent human validation of the factuality labels was collected in this
study, and it is worth stating precisely what this does and does not leave open.
The repeated runs in Chapter~\ref{ch:results},
Section~\ref{sec:results-final-campaign} establish that the difference between the
two preference sets is repeatedly observed: it survives changes in fold assignment
and training stochasticity, so it is not confined to a single favourable run.

Reproducibility and correctness are separate properties, however. That a
preference relation is consistently recoverable says nothing about whether the
underlying AlignScore judgments agree with what a human reader would consider
faithful to the source. A judge can produce labels that are highly reproducible
and systematically wrong, and nothing in this study would detect that case. The
labels function here as the reference standard, not as verified ground truth, and
every result is conditional on them in that sense. Establishing agreement with
human factuality judgments requires annotation this study did not collect,
quantified against a chance-corrected agreement measure such as Cohen's $\kappa$
\citep{cohen1960coefficient}, and the question is recorded as an open limitation
in Chapter~\ref{ch:conclusion}. The concern is not peculiar to this study:
reviews of reinforcement learning from human feedback identify the quality of the
preference data itself, and the difficulty of auditing it, as a central open
problem \citep{casper2023open}, and substituting an automatic judge for human
annotators relocates that problem rather than removing it.

\section{Stochastic Variation Between Runs}
\label{sec:disc-stochasticity}

Run-to-run variation is large enough in this setting to determine what conclusions
can be drawn, and it deserves treatment as a finding rather than as an
inconvenience. Two same-configuration runs differed by 4.7 percentage points
($+0.060$ versus $+0.013$), and across the six runs the gap ranged from $-0.010$
to $+0.060$. A study reporting only one of these runs could honestly have
concluded either that GCA helps substantially or that it does not help at all.

This is a known hazard rather than a peculiarity of the present setup. Fine-tuning
pretrained language models is sensitive to random seed alone, with weight
initialization and data order contributing comparably to the variance of
out-of-sample performance \citep{dodge2020finetuning}, and reporting score
distributions across several runs rather than a single figure has been urged for
precisely this reason \citep{reimers2017reporting}. The repeated-run design used
here follows that recommendation, and the spread it exposes is the reason this
thesis reports a distribution of gaps rather than a single headline number.

This variation traced back to an implementation property documented in
Chapter~\ref{ch:implementation}, Section~\ref{sec:impl-rm-training}: in the
original six-run campaign, the seed argument fixed the cross-validation fold
assignment through Python's standard random number generator, but PyTorch's
generator was not seeded in the cross-validation path. Several sources of
training stochasticity were therefore uncontrolled, including reward-head weight
initialization, the order in which training batches are drawn, and dropout
masks. Runs were consequently not exactly reproducible from a seed value alone.
This had an interpretive consequence at the time: the two executions at seed 42
were genuine repetitions rather than duplicates, because they shared the fold
assignment but not the training stochasticity, and the 4.7-point spread between
them reflects that uncontrolled stochasticity rather than fold assignment alone.

This property has since been fixed --- PyTorch and CUDA are now seeded per fold,
with an explicit generator for the training \texttt{DataLoader} and seeded
worker initialization, so a run is fully determined by its seed on fixed
hardware. The extended twenty-run campaign in Chapter~\ref{ch:results},
Section~\ref{sec:results-extended-campaign} was carried out under this corrected
procedure; each of its twenty seeds is consequently a genuinely independent and
exactly reproducible sample, which is what made it meaningful to extend the
campaign at all. The fix does not retroactively change the original six runs,
which remain a documented instance of exactly the seed-sensitivity the
literature describes, but it does mean the concern raised in this section no
longer applies to the results this thesis reports as its headline finding.

The practical implication generalizes beyond this study regardless of the fix.
At the accuracy levels and dataset sizes involved here, differences of one to two
percentage points between conditions are within the range that a single
under-seeded run can produce by chance. Comparisons of preference-construction
methods reported from single training runs at this scale, without controlling
every source of training randomness, should be treated with corresponding
caution.

\section{Confounds That Remain}
\label{sec:disc-confounds}

Two confounds could not be removed within this design, and both plausibly affect
the reported numbers.

The candidate pairs differ in decoding temperature, and
Chapter~\ref{ch:methodology}, Table~\ref{tab:temperature-asymmetry} shows the
consequence is not small: the lower-temperature candidate is preferred in 68.5\%
of holistic pairs and 63.0\% of GCA pairs on the data where this could be
measured. Because temperature systematically affects surface properties such as
length and lexical conservativeness, a reward model could recover a meaningful
share of the preference relation from stylistic cues without representing
factuality at all. The asymmetry also differs between conditions by roughly five
percentage points, so the two preference sets are not equivalent with respect to
this confound, and part of the observed difference between them could in principle
reflect that rather than granularity. Sampling both candidates at the same
temperature would remove the confound and is the cleanest correction available.

The reward model also sees only a truncated prefix of each source article
(Chapter~\ref{ch:implementation}, Section~\ref{sec:impl-rm-training}), while
subset selection admits articles several times longer than the retained prefix.
Where the evidence bearing on a claim lies beyond that prefix, the corresponding
preference is not recoverable from the model's input even in principle. This is a
plausible contributor to the low absolute accuracies discussed in
Section~\ref{sec:disc-weak-signal}, though its magnitude was not measured, and it
should not be assumed to be the dominant cause without the ablation described
there.

\section{Generalizability}
\label{sec:disc-generalizability}

Every result reported here is obtained under one combination of choices:
CNN/DailyMail as the corpus, AlignScore as the evaluator,
Mistral-7B-Instruct-v0.3 as the candidate generator, and a RoBERTa-base
Bradley--Terry model as the reward model. Each of these bounds the conclusions in
a distinct way.

The evaluator is the most consequential. Section~\ref{sec:disc-interpretation}
raises the possibility that whether external decomposition helps is related to
whether the evaluator already decomposes internally, which would make the finding
partly \emph{about} the evaluator's architecture rather than simply conditional on
AlignScore as a fixed instrument. If that reading is correct, a different
factuality model with different internal aggregation would shift the result. The corpus matters because CNN/DailyMail
summaries are short, extractive-leaning, and drawn from a single register; domains
with longer or more abstractive summaries would change both the number of
sentences per summary and the difficulty of the underlying factuality judgments.
The reward-model backbone matters because a stronger encoder, or one with a longer
context window, might recover more of the preference relation and thereby narrow
or widen the difference between conditions.

Taken together, these dependencies mean the study is best read as characterizing
behaviour under one specific configuration, and as motivating a hypothesis about
the interaction between external and evaluator-internal decomposition, rather than
as establishing a quantity that transfers directly to other setups.

\chapter{Conclusion}
\label{ch:conclusion}

\section{Summary}
\label{sec:concl-summary}

This thesis asked a controlled question: does converting sentence-level factuality
scores into a summary-level preference through Granular Credit Assignment produce a
more learnable reward-model signal than scoring each candidate summary
holistically. The two conditions were compared under an identical candidate pool,
an identical factuality judge, and an identical Bradley--Terry training recipe, so
that the only manipulated factor was the preference-construction procedure itself.

Averaged over six repeated runs at $n=1{,}000$, GCA preferences yield an
advantage of approximately $+3.08$ percentage points in reward-model pairwise
accuracy. The advantage does not remain equally large as the dataset grows: it is
approximately $-0.42$ percentage points at $n=5{,}000$ and $+0.35$ percentage
points at $n=10{,}000$, while absolute accuracy improves for both conditions.
Uncertainty estimates computed over cross-validation folds are anti-conservative,
because folds within a run are not independent, so the direction and approximate
magnitude of the effect are better supported than any exact significance level.

The experiments therefore do not support a universal claim that GCA is preferable
to holistic factuality supervision. They show that sentence-level aggregation can
improve the learnability of automatically generated pairwise preferences under
particular configurations, most clearly in the $n=1{,}000$ setting. In the
exploratory mode comparison, the advantage appears under evaluator modes that do
not themselves decompose the input, which motivates the hypothesis that external
and evaluator-internal decomposition may act partly as substitutes. At larger
sample sizes the observed advantage is small and inconsistent in sign. Taken
together, the results indicate that the usefulness of granular credit assignment
varies with evaluator configuration, dataset size, and aggregation choice rather
than following from finer granularity as such.

\section{Answers to the Research Questions}
\label{sec:concl-rqs}

\textbf{RQ1} (does GCA produce a more learnable preference signal than holistic
scoring?) is answered \emph{conditionally}. An average advantage of approximately
3 percentage points is repeatedly observed at $n=1{,}000$, favouring GCA in five
of six runs, though the run-level test does not meet the conventional threshold.
It is not observed in the single run at $n=5{,}000$ and is marginal at
$n=10{,}000$. The answer is therefore configuration-dependent rather than
general.

\textbf{RQ2} (how do evaluator configuration and aggregation strategy affect
relative learnability?) is answered \emph{associationally}. Both the aggregation
formula and the evaluator mode are associated with substantial changes in the
outcome, and in the exploratory sweep the two modes in which the evaluator
performs its own internal decomposition are also the two in which GCA shows no
advantage (Chapter~\ref{ch:results}, Section~\ref{sec:results-mode-sweep}).
Because only one mode received repeated-run validation, this pattern is suggestive
rather than established.

\textbf{RQ3} (how stable are the differences across runs and dataset sizes?) is
answered \emph{partially}. The difference is directionally consistent across six
runs at fixed $n=1{,}000$, favouring GCA in five of six. It is not reproduced with
comparable magnitude in the single runs at $n=5{,}000$ and $n=10{,}000$, which use
nested supersets of the same data (Section~\ref{sec:disc-scale}).

\section{Limitations}
\label{sec:concl-limitations}

\begin{enumerate}
  \item \textbf{No human factuality audit.} No independent human annotation was
    collected in the final study, so the evaluator's judgments were never checked
    against human assessments of faithfulness
    (Chapter~\ref{ch:discussion}, Section~\ref{sec:disc-reliability}).
  \item \textbf{The labels are not independent ground truth.} AlignScore supplies
    both the holistic and the sentence-level scores, so it functions as the
    reference standard rather than as a verified one. Every result is conditional
    on that evaluator.
  \item \textbf{Preference learnability is not factual correctness.} The
    dependent variable measures how consistently a preference relation can be
    recovered by the chosen architecture, not whether the preferences are
    factually right (Chapter~\ref{ch:methodology},
    Section~\ref{sec:methodology-scope}).
  \item \textbf{No downstream policy evaluation.} The final study contains no
    DPO or other policy-optimization stage, so it does not establish whether
    models trained on either preference set generate more factually consistent
    summaries.
  \item \textbf{Both reward models are weak in absolute terms.} At $n=1{,}000$
    the conditions reach 0.5295 and 0.5603 pairwise accuracy against a chance
    level of 0.5, so the comparison is between two weak signals
    (Chapter~\ref{ch:discussion}, Section~\ref{sec:disc-weak-signal}).
  \item \textbf{Cross-validation folds are not independent.} Uncertainty
    estimates computed over 30 folds treat correlated observations as
    independent and are therefore anti-conservative; the direction and
    approximate magnitude of the effect are better supported than any exact
    significance level (Section~\ref{sec:disc-stats-caveat}).
  \item \textbf{Incomplete determinism across runs.} Subset selection and
    cross-validation fold assignment are explicitly seeded, but PyTorch's
    generator is not seeded in the cross-validation path, leaving weight
    initialization, batch ordering and dropout uncontrolled
    (Chapter~\ref{ch:implementation}, Section~\ref{sec:impl-rm-training}). Exact
    run-level reproduction from the command-line seed alone is therefore not
    guaranteed.
  \item \textbf{Configuration selected on observed results.} The aggregation
    strategy and evaluator mode were chosen after inspecting $n=1{,}000$ outcomes
    (Chapter~\ref{ch:methodology}, Section~\ref{sec:methodology-exploratory}), and
    the larger-scale subsets are nested supersets of that same data
    (Section~\ref{sec:methodology-nesting}), so no result in this thesis
    evaluates the selected configuration on independent data.
  \item \textbf{Candidate generation confounds temperature with condition.} The
    two candidates differ in decoding temperature, and the lower-temperature
    candidate is preferred substantially more often under both conditions and to
    differing degrees (Chapter~\ref{ch:methodology},
    Section~\ref{sec:methodology-ethics}), so temperature-correlated features may
    carry part of the label information.
  \item \textbf{Source-document truncation.} The reward model receives only a
    truncated prefix of each article, so evidence located later in the source is
    unavailable to it (Chapter~\ref{ch:implementation},
    Section~\ref{sec:impl-rm-training}).
  \item \textbf{Evaluator and dataset specificity.} All findings are obtained on
    CNN/DailyMail with AlignScore, Mistral-7B-Instruct-v0.3 candidates, and a
    RoBERTa-base reward model, and need not transfer to other domains,
    evaluators, or architectures.
  \item \textbf{Sentence segmentation was not evaluated.} Summaries are split
    with a regex-based splitter whose accuracy was never measured
    (Chapter~\ref{ch:implementation}, Section~\ref{sec:impl-preferences}).
  \item \textbf{Nested dataset sizes.} The $n=1{,}000$, $n=5{,}000$ and
    $n=10{,}000$ subsets overlap almost completely
    (Chapter~\ref{ch:methodology}, Section~\ref{sec:methodology-nesting}), so the
    larger-scale runs are not independent samples and no experiment evaluates the
    selected configuration on held-out data.
  \item \textbf{No error-category analysis.} No analysis links the observed
    effects to specific categories of factual error; this was not part of the
    final research questions and is listed under future work.
\end{enumerate}

\section{Future Work}
\label{sec:concl-future}

\begin{itemize}
  \item \textbf{Repeated runs at larger scale.} The $n=5{,}000$ and $n=10{,}000$
    experiments were each executed once. Repeating them would test whether their
    intervals overlap the $n=1{,}000$ estimate, which would indicate run-to-run
    variance, or remain separated, which would support a dataset-size effect
    (Chapter~\ref{ch:discussion}, Section~\ref{sec:disc-scale}).
  \item \textbf{Remove the temperature confound.} Generating both candidates at a
    single decoding temperature, and comparing the resulting preference sets
    against the current ones, would establish how much of the observed structure
    is attributable to temperature-correlated surface features rather than to
    factuality (Section~\ref{sec:disc-confounds}).
  \item \textbf{Measure the effect of source truncation.} Varying the retained
    article length while holding all else fixed would quantify how much of the
    near-chance reward-model accuracy is attributable to evidence falling outside
    the model's input window (Section~\ref{sec:disc-weak-signal}).
  \item \textbf{Test the decomposition-interaction mechanism directly.}
    Section~\ref{sec:disc-interpretation} infers from AlignScore's documented
    behavior that its split modes perform internally much of what GCA performs
    externally. Comparing per-sentence scores under split and unsplit modes on the
    same pairs would test that inference rather than rest on it.
  \item \textbf{Independent validation of the labels.} Human annotation of a
    stratified sample of holistic/GCA disagreement cases, or scoring the same
    candidate pool with a second factuality model, would provide the external
    reference this study lacks and would indicate whether either construction
    procedure produces labels closer to human judgment.
  \item \textbf{Downstream policy evaluation.} Training a summarization policy on
    each preference set and evaluating the factual consistency of its generations
    would test whether differences in preference learnability translate into
    differences in generated output, which this study deliberately does not
    address (Chapter~\ref{ch:methodology}, Section~\ref{sec:methodology-evolution}).
  \item \textbf{Error-category analysis.} Categorizing disagreement cases by
    error type would test whether any advantage concentrates in localized errors
    such as entity or relation mistakes, as originally hypothesized.
\end{itemize}


\bibliographystyle{plainnat}
\bibliography{references}

\appendix
\chapter{Data and Code Availability}
\label{app:availability}

All source code, configuration files, Slurm job scripts, and analysis notebooks
used to produce the results reported in this thesis are publicly available at:

\begin{quote}
\url{https://github.com/hasnat23/master-thesis-rlaif-gca}
\end{quote}

The repository is organized as follows. \url{src/} contains the pipeline
implementation, split into the modules referenced throughout
Chapter~\ref{ch:implementation}: \url{src/data/} (subset selection and record
schemas), \url{src/generation/} (candidate summary generation),
\url{src/judging/} (the AlignScore judge wrapper and both preference-construction
regimes), \url{src/reward_model/} (Bradley--Terry training and cross-validation),
and \url{src/eval/} and \url{src/analysis/} (evaluation and statistics).
\url{configs/} holds the YAML configuration files for subset selection,
generation, and judging. \url{slurm/} holds the batch scripts submitted on MOGON
NHR, including \url{slurm/train_rm_1000.sh}, which produced the headline
$n=1{,}000$ result. \url{analysis/} holds the standalone scripts used for the
formula sweep, the judge-mode sweep, and the holistic-versus-GCA comparison,
together with \url{analysis/aggregate_campaigns.py}, which computes the
run-level means, bootstrap intervals and Wilcoxon tests for each seed campaign,
and \url{analysis/effect_sizes_and_equivalence.py}, which computes the effect
sizes, equivalence tests and minimum detectable effects reported in
Chapter~\ref{ch:results}, Section~\ref{sec:results-effect-sizes}. Both read the
per-run JSON summaries under \url{outputs/} directly, so every reported figure
can be regenerated from the stored run artifacts without retraining.
\url{progress-updates/} contains the dated working notes kept over the course of
the project.

The underlying corpus, CNN/DailyMail version \texttt{3.0.0}
\citep{hermann2015teaching, nallapati2016abstractive}, is publicly available
through the Hugging Face \texttt{datasets} library and is not redistributed in the
repository. The candidate summaries, preference files, and trained reward-model
checkpoints generated during this project are likewise not committed, as they run
to several gigabytes; they can be regenerated from the corpus using the
configuration recorded in Appendix~\ref{app:reproduction}. The factuality
evaluator, \texttt{yzha/AlignScore} \citep{zha2023alignscore}, and the candidate
generator, Mistral-7B-Instruct-v0.3, are both publicly released model checkpoints
obtained from the Hugging Face Hub.

\chapter{Reproduction Configuration}
\label{app:reproduction}

This appendix records the exact configuration required to reproduce the headline
$n=1{,}000$ result. It restates in executable form what
Chapter~\ref{ch:methodology}, Table~\ref{tab:methodology-summary} summarizes and
Chapter~\ref{ch:implementation} describes in prose. Under the corrected,
fully-seeded training procedure described in
Chapter~\ref{ch:implementation}, Section~\ref{sec:impl-rm-training}, a single
run of these commands at a given seed value is deterministic on fixed
hardware; running the full twenty-run headline campaign means repeating the
reward-model training command below once per seed in $\{1, \ldots, 20\}$, as
\url{slurm/seed_campaign.sh} does.

\section{Summarization Prompt}
\label{app:repro-prompt}

Both candidate summaries for a given article are generated from the same fixed
prompt template (\texttt{SUMMARIZATION\_PROMPT} in
\url{src/generation/candidates.py}); the two candidates differ only in the
sampling parameters applied at decoding time.

\begin{quote}
\small\ttfamily
Summarize the following news article in a concise paragraph.\\[2mm]
Article:\\
\{article\}\\[2mm]
Summary:
\end{quote}

\section{Pipeline Parameters}
\label{app:repro-params}

Table~\ref{tab:app-params} lists the parameters of each pipeline stage as used for
the final comparison. Where a value differs from the default recorded in the
corresponding YAML file, the command-line override in
Section~\ref{app:repro-commands} is authoritative; this is the case for the
subset size, the tie margin, and the GCA exponent $\alpha$, each of which changed
over the course of the project as described in Chapter~\ref{ch:methodology},
Section~\ref{sec:methodology-evolution}.

\begin{table}[htbp]
\centering
\caption[Pipeline parameters for the final $n=1{,}000$ comparison]{Pipeline
parameters for the final $n=1{,}000$ comparison.}
\label{tab:app-params}
\begin{tabularx}{\textwidth}{llX}
\toprule
Stage & Parameter & Value \\
\midrule
\multirow{5}{*}{Subset selection}
  & Corpus              & CNN/DailyMail 3.0.0, \texttt{test} split \\
  & Samples             & 1{,}000 \\
  & Subset seed         & 200 \\
  & Max article length  & 8{,}000 characters \\
  & Max summary length  & 2{,}000 characters \\
\midrule
\multirow{6}{*}{Candidate generation}
  & Model               & Mistral-7B-Instruct-v0.3, \texttt{bfloat16} \\
  & Candidate A         & $T=0.7$, top-$p=0.9$ \\
  & Candidate B         & $T=1.0$, top-$p=0.95$ \\
  & Max input tokens    & 1{,}024 \\
  & Max new tokens      & 256 \\
  & Batch size          & 4 \\
\midrule
\multirow{4}{*}{Preference construction}
  & Evaluator           & \texttt{yzha/AlignScore}, \texttt{FacebookAI/roberta-base} \\
  & Evaluation mode     & \texttt{nli} \\
  & GCA exponent        & $\alpha = 0.0$ (simple mean) \\
  & Tie margin          & $0$ \\
\midrule
\multirow{7}{*}{Reward-model training}
  & Backbone            & \texttt{FacebookAI/roberta-base} \\
  & Epochs              & 5 \\
  & Batch size          & 8 \\
  & Learning rate       & $2\times10^{-5}$, linear warmup over 10\% of steps \\
  & Max sequence length & 512 tokens \\
  & Max article chars   & 2{,}000 \\
  & Cross-validation    & 5-fold; training seeds 1--20 (headline campaign,
    Chapter~\ref{ch:results}, Table~\ref{tab:seed-campaign-full}) \\
\bottomrule
\end{tabularx}
\end{table}

\section{Command-Line Invocations}
\label{app:repro-commands}

The three stages are run in sequence. Preference construction produces both
conditions from a single invocation, and reward-model training consumes both
preference files so that the holistic and GCA models are trained under an
identical recipe within one process.

\begin{quote}
\footnotesize
\begin{verbatim}
# 1. Subset selection
python scripts/01_prepare_subset.py --n-samples 1000 --seed 200

# 2. Candidate generation
python scripts/02_generate_candidates.py \
    --config configs/generation_1000.yaml

# 3. Preference construction (both conditions)
python src/judging/build_reward_preferences.py \
    --mode both --alpha 0.0 --margin 0 \
    --alignscore-evaluation-mode nli \
    --alignscore-backbone FacebookAI/roberta-base

# 4. Reward-model training (5-fold CV, both conditions)
python src/reward_model/run_training.py \
    --holistic data/preferences_1000/\
holistic_reward_preferences_1000.jsonl \
    --gca      data/preferences_1000/\
gca_reward_preferences_1000.jsonl \
    --output-dir outputs/reward_models_1000 \
    --backbone FacebookAI/roberta-base \
    --epochs 5 --batch-size 8 --lr 2e-5 \
    --max-length 512 --max-article-chars 2000 \
    --kfold 5 --seed 1
\end{verbatim}
\end{quote}

\noindent
This reproduces one run of the headline campaign, seed 1 of 20. Substituting
each of \texttt{--seed 2} through \texttt{--seed 20} in turn reproduces the
remaining nineteen.

\noindent
In step 4 the two preference-file paths are shown split across lines for
typesetting only; the trailing backslash inside a path is a line continuation in
this listing, not part of the filename.

On the MOGON NHR cluster these are submitted as Slurm batch jobs
(\url{slurm/generate_candidates_1000.sh}, \url{slurm/build_preferences_1000.sh},
\url{slurm/train_rm_1000.sh}) on partition \texttt{a100dl} with one
A100-SXM4-40GB GPU, 32\,GB of host memory, and four CPU cores per task.

\chapter{Example Preference Record}
\label{app:example-record}

Every record written by the GCA preference-construction stage retains the full
sentence-level detail behind the summary-level decision, which is what makes the
qualitative disagreement analysis in Chapter~\ref{ch:results} possible. The
listing below shows one such record, abridged: the article, reference summary, and the
\texttt{chosen}/\texttt{rejected} fields are truncated, and the sentence list for
summary B is cut after the first three sentences. All numeric values are
reproduced verbatim.

The record shown was produced by the earlier $n=200$ run, which used the
exploratory settings $\alpha = 0.5$ and margin $=0.05$ rather than the final
$\alpha = 0.0$ and margin $=0$ (Chapter~\ref{ch:methodology},
Section~\ref{sec:methodology-evolution}). It is included to illustrate the record
\emph{schema} and the relationship between the sentence scores and the aggregate,
not as an instance of the final configuration. Note in particular that
\texttt{gca\_score\_a} $=0.382$ is well below \texttt{mean\_score\_a} $=0.629$:
this is the penalty term at $\alpha = 0.5$ pulling the aggregate down towards the
summary's weakest sentence, the behaviour that Chapter~\ref{ch:results},
Section~\ref{sec:results-formula-opt} found to perform worse than the simple mean.

\begin{quote}
\footnotesize
\begin{verbatim}
{
  "sample_id": "cnn_00000_b0b34b2d843e",
  "article": "LONDON, England (Reuters) -- Harry Potter star Daniel ...",
  "reference_summary": "Harry Potter star Daniel Radcliffe gets ...",
  "summary_a": "Daniel Radcliffe, the 18-year-old actor who plays ...",
  "summary_b": "On Monday, July 23, 2007, Harry Potter actor Daniel ...",
  "gca_score_a": 0.382102,
  "gca_score_b": 0.169184,
  "gca_score_diff": 0.212918,
  "mean_score_a": 0.629039,
  "mean_score_b": 0.578186,
  "min_score_a": 0.232103,
  "min_score_b": 0.049505,
  "n_sentences_a": 3,
  "n_sentences_b": 6,
  "n_low_faithfulness_a": 1,
  "n_low_faithfulness_b": 3,
  "decision": "A",
  "judge_method": "gca",
  "judge_model": "yzha/AlignScore [FacebookAI/roberta-base]",
  "margin": 0.05,
  "alpha": 0.5,
  "sentence_details_a": [
    { "index": 0,
      "text": "Daniel Radcliffe, the 18-year-old actor who plays Harry
               Potter, has gained access to a 20 million pound fortune
               on his 18th birthday.",
      "score": 0.748840, "label": "faithful" },
    { "index": 1,
      "text": "Despite the newfound wealth, he has stated that he does
               not plan to be extravagant, and will instead use his
               money to buy books, CDs, and DVDs.",
      "score": 0.232103, "label": "hallucinated" },
    { "index": 2,
      "text": "Radcliffe has filmed a TV movie and an Australian film,
               and has made his stage debut in \"Equus.\" He is expected
               to face closer media scrutiny now that he is legally an
               adult.",
      "score": 0.906173, "label": "faithful" }
  ],
  "sentence_details_b": [
    { "index": 0, "text": "On Monday, July 23, 2007, ...",
      "score": 0.475803, "label": "hallucinated" },
    { "index": 1, "text": "Radcliffe has stated he has no ...",
      "score": 0.997014, "label": "faithful" },
    { "index": 2, "text": "Instead, he plans to use the money ...",
      "score": 0.049505, "label": "hallucinated" }
  ]
}
\end{verbatim}
\end{quote}

Two features of this record are worth drawing out, because they illustrate
mechanisms discussed in the main text. First, the \texttt{label} field is a
threshold applied at a score of 0.5 for reporting convenience only; it plays no
part in the aggregation, which operates on the continuous scores. Second, summary
B contains both the highest-scoring sentence in the pair (0.997) and the
lowest (0.050). Under the holistic condition this internal variation is not
observable at all: the summary receives a single score. This is precisely the
information that sentence-level decomposition exposes and holistic scoring
discards, and whether exposing it yields a more learnable preference signal is the
question the thesis sets out to answer.


\end{document}
